\documentclass[../../main.tex]{subfiles}
\graphicspath{{\subfix{../../image/}}}

\begin{document}

\section{差商与扩张定理}\label{section:ch05-延拓}

\subsection{差商与 $W^{1,p}$}\label{subsubsection:ch05-差商W1p}

\begin{definition}\label{definition:ch05-差商}
设$U\subseteq \mathbb{R}^n$是开集,$u:U\to\mathbb{R}$ 为局部可和函数,即$u\in L_{\text{loc}}^1(U)$,且$V\subset \subset U$。
步长为 $h$ 的第 $i$ 差商定义为
\begin{align*}
D_{i}^{h}u(x)\triangleq \frac{u(x+he_{i})-u(x)}{h}\qquad (i=1,\cdots,n),
\end{align*}
其中 $x\in V$，$h\in\mathbb{R}$，$0<|h|<\operatorname{dist}(V,\partial U)$。
并且$D^{h}u\triangleq(D_{1}^{h}u,\cdots,D_{n}^{h}u)$。
\end{definition}
\begin{remark}
差商来自导数定义中的有限增量比. 对于通常的偏导数, 我们考察
\begin{align*}
\frac{\partial u}{\partial x_i}(x)
=
\lim\limits_{h\to0}
\frac{u(x+he_i)-u(x)}{h},
\end{align*}
前提是这个极限存在. 差商则保留极限之前的表达式, 即先固定一个非零步长 $h$, 定义
\begin{align*}
D_i^h u(x)
=
\frac{u(x+he_i)-u(x)}{h}.
\end{align*}
因此, 差商可以理解为沿第 $i$ 个坐标方向、在尺度 $|h|$ 上的平均变化率.
定义差商时不需要预先知道 $u$ 可微, 这使它能够用于研究尚未确定具有弱导数的函数.

分母中的 $h$ 起到按位移大小归一化的作用. 若 $u$ 在 $x$ 附近可微, 则
\begin{align*}
u(x+he_i)-u(x)
=
h\frac{\partial u}{\partial x_i}(x)+o(|h|),
\end{align*}
所以除以 $h$ 后, 主项恰好是偏导数. 如果只考虑函数值之差而不除以 $h$, 则对于连续函数, 这个差通常直接趋于零, 无法反映导数的大小.

这里 $e_i$ 是第 $i$ 个标准单位向量, 因而 $x+he_i$ 只改变 $x$ 的第 $i$ 个坐标. 将所有坐标方向的差商放在一起, 得到
\begin{align*}
D^h u
\triangleq
(D_1^h u,\cdots,D_n^h u),
\end{align*}
它是梯度的离散近似. 上标 $h$ 表示步长参数, 并不表示求 $h$ 阶导数.
允许 $h$ 取正值或负值, 是为了同时考察两个方向的增量. 当 $h<0$ 时, 分母也带有负号, 所以两种方向的差商在极限存在时都趋于同一个偏导数.

条件 $U$ 为开集, 保证每个内点附近都有一小段可以移动而不离开定义域的空间.
但是, 如果让 $x$ 遍历整个 $U$, 则靠近边界时, 一般无法找到一个对所有 $x$ 都适用的非零步长.
因此, 我们先在一个远离边界的局部区域 $V$ 上讨论差商.

由\subcref{Real Analysis-proposition:点集间的距离的性质}{Real Analysis-proposition:点集间的距离的性质-3}, 有
\begin{align*}
d
\triangleq
\operatorname{dist}(V,\partial U)
=
\operatorname{dist}(\overline{V},\mathbb{R}^n\setminus U)
>0.
\end{align*}
因此, 当 $0<|h|<d$ 时, 对任意 $x\in\overline{V}$ 和 $t\in[0,1]$, 都有
\begin{align*}
|x+the_i-x|
=
t|h|
<d,
\end{align*}
从而
\begin{align*}
x+the_i\in U.
\end{align*}
这不仅保证差商中的两个取值点 $x$ 和 $x+he_i$ 都属于 $U$, 还保证连接它们的整条线段都在 $U$ 内.
严格不等式 $|h|<d$ 留出了正的边界距离, 使平移后的闭包仍然紧含于 $U$.
条件 $h\neq0$ 则是因为差商的分母不能为零.
如果 $U=\mathbb{R}^n$, 约定到空集的距离为 $+\infty$, 此时没有来自边界的步长限制.

要求 $u\in L_{\mathrm{loc}}^1(U)$, 是为了使差商适用于弱导数的积分框架.
差商的点值公式本身对一般函数也可以写出, 但后续需要将差商与测试函数相乘并积分, 因而需要局部可积性.
对于每个固定的允许步长 $h$, 由于 $V$ 及其平移 $V+he_i$ 都相对紧于 $U$, 有
\begin{align*}
\int_V |D_i^h u(x)|\,\mathrm{d}x
&\leqslant
\frac{1}{|h|}
\left(
\int_V |u(x+he_i)|\,\mathrm{d}x
+
\int_V |u(x)|\,\mathrm{d}x
\right)\\
&=
\frac{1}{|h|}
\left(
\int_{V+he_i}|u(y)|\,\mathrm{d}y
+
\int_V |u(x)|\,\mathrm{d}x
\right)
<\infty.
\end{align*}
因此 $D_i^h u\in L^1(V)$.
这里不必要求 $u$ 在整个 $U$ 上可积, 因为计算只涉及 $U$ 内部的相对紧区域.

差商与导数之间最直接的联系来自微积分基本定理.
若 $u\in C^1(U)$, 则由上述线段包含关系, 对任意允许的 $x$ 和 $h$, 有
\begin{align*}
u(x+he_i)-u(x)
=
\int_0^1
\frac{\mathrm{d}}{\mathrm{d}t}u(x+the_i)\,\mathrm{d}t=
h\int_0^1
\frac{\partial u}{\partial x_i}(x+the_i)\,\mathrm{d}t.
\end{align*}
所以
\begin{align*}
D_i^h u(x)
=
\int_0^1
\frac{\partial u}{\partial x_i}(x+the_i)\,\mathrm{d}t.
\end{align*}
这说明, 对光滑函数而言, 差商正是偏导数沿一小段线段的平均值.
当 $h\to0$ 时, 这段线段收缩到 $x$, 差商便趋于该点的偏导数.
这也是差商能够近似导数、并能够由导数估计控制的原因.

对于局部可积函数, 函数值及差商都应按几乎处处的意义理解.
如果改变 $u$ 在某个零测集上的取值, 则对于每个固定的 $h$, 差商只可能在该零测集及其平移上发生变化.
由于平移保持零测性, 差商所确定的可积函数等价类不受代表元选择的影响.现在严格证明差商不受代表元选择的影响.
\begin{proof}
设 $u$ 和 $\widetilde{u}$ 是两个代表元, 并且存在零测集 $N\subseteq U$, 使得
\begin{align*}
u(y)=\widetilde{u}(y),
\quad y\in U\setminus N.
\end{align*}
固定一个允许的非零步长 $h$满足:$h\in\mathbb{R}$，$0<|h|<\operatorname{dist}(V,\partial U)$.
比较两个差商时, 要同时比较两组函数值:
\begin{align*}
D_i^h u(x)
=
\frac{u(x+he_i)-u(x)}{h},\quad 
D_i^h\widetilde{u}(x)
=
\frac{\widetilde{u}(x+he_i)-\widetilde{u}(x)}{h}.
\end{align*}
只有当 $x$ 或 $x+he_i$ 落入 $N$ 时, 这两个差商才可能不同.
其中, 条件 $x+he_i\in N$ 等价于
\begin{align*}
x\in N-he_i.
\end{align*}
因此, 两个差商可能不同的点只会出现在集合
\begin{align*}
E_h\triangleq V\cap\bigl(N\cup(N-he_i)\bigr)
\end{align*}
中.
由\hyperref[Real Analysis-theorem:外测度的平移不变性]{测度的平移不变性}以及有限个零测集之并仍为零测集, 有
\begin{align*}
|E_h|
\leqslant
|N|+|N-he_i|
=0.
\end{align*}
对于每个 $x\in V\setminus E_h$, 两个取值点都不属于 $N$, 所以
\begin{align*}
u(x)=\widetilde{u}(x),\quad
u(x+he_i)=\widetilde{u}(x+he_i),
\end{align*}
从而
\begin{align*}
D_i^h u(x)=D_i^h\widetilde{u}(x).
\end{align*}
这就证明了两个差商在 $V$ 上几乎处处相等, 因而代表同一个可积函数等价类.
\end{proof}
\end{remark}

\begin{proposition}[紧集平移的紧性]\label{proposition:紧集平移与紧参数平移集}
设 $K,A\subseteq\mathbb{R}^n$ 都是紧集, 则
\begin{align*}
K+A
\triangleq
\{x+a:x\in K,\, a\in A\}
=
\bigcup\limits_{a\in A}(K+a)
\end{align*}
也是紧集.
特别地, 有以下结论:
\begin{enumerate}[(1)]
\item 对任意固定的向量 $a\in\mathbb{R}^n$, 平移集
\begin{align*}
K+a\triangleq\{x+a:x\in K\}
\end{align*}
是紧集.
\item 对任意有限个向量 $a_1,\cdots,a_m\in\mathbb{R}^n$, 集合
\begin{align*}
\bigcup\limits_{j=1}^m
\{x+ta_j:x\in K,\, t\in[0,1]\}
\end{align*}
是紧集.
\end{enumerate}
\end{proposition}
\begin{proof}
由于 $K$ 和 $A$ 都是紧集, 它们的乘积 $K\times A$ 是 $\mathbb{R}^{2n}$ 中的紧集.
定义映射
\begin{align*}
F:K\times A\longrightarrow\mathbb{R}^n,\quad
(x,a)\longmapsto x+a.
\end{align*}
对任意 $(x,a),(y,b)\in K\times A$, 有
\begin{align*}
|F(x,a)-F(y,b)|
\leqslant |x-y|+|a-b|\leqslant
\sqrt{2}\sqrt{|x-y|^2+|a-b|^2}.
\end{align*}
因此 $F$ 连续.
\hyperref[Topology-theorem:topology-紧集的连续像是紧集]{紧集在连续映射下的像仍然是紧集}, 故
\begin{align*}
K+A=F(K\times A)
\end{align*}
是紧集.

对于结论(1), 取紧集 $A=\{a\}$, 即可得到 $K+a$ 是紧集.

对于结论(2), 令
\begin{align*}
A_j\triangleq\{ta_j:t\in[0,1]\},
\quad j=1,\cdots,m.
\end{align*}
映射 $t\mapsto ta_j$ 连续, 所以 $A_j$ 是紧区间 $[0,1]$ 的连续像, 从而由\hyperref[Topology-theorem:topology-紧集的连续像是紧集]{紧集在连续映射下的像仍然是紧集}知$A_j$都是紧集.
因此, 有限并
\begin{align*}
A\triangleq\bigcup\limits_{j=1}^m A_j
\end{align*}
也是紧集.
应用已经证明的结论, 可知
\begin{align*}
K+A
=
\bigcup\limits_{j=1}^m
\{x+ta_j:x\in K,\, t\in[0,1]\}
\end{align*}
是紧集.
\end{proof}

\begin{proposition}[开集内部小幅平移后仍在原本开集中]\label{proposition:开集内部小幅平移后仍在原本开集中}
设 $U\subseteq\mathbb{R}^n$ 是开集, $\varnothing\neq V\subset\subset U$, 记
\begin{align*}
\delta\triangleq\operatorname{dist}(V,\partial U)>0.
\end{align*}
当 $U=\mathbb{R}^n$ 时, 约定 $\delta=+\infty$.
则对任意满足 $|a|<\delta$ 的向量 $a\in\mathbb{R}^n$, 都有
\begin{align*}
\{x+ta:x\in\overline{V},\, t\in[0,1]\}
\subset\subset U.
\end{align*}
特别地, $\overline{V}+a\subset\subset U$.

更一般地, 若有限个向量 $a_1,\cdots,a_m$ 都满足 $|a_j|<\delta$, 则
\begin{align*}
\bigcup\limits_{j=1}^m
\{x+ta_j:x\in\overline{V},\, t\in[0,1]\}
\subset\subset U.
\end{align*}
\end{proposition}
\begin{proof}
当 $U=\mathbb{R}^n$ 时, 结论直接由\cref{proposition:紧集平移与紧参数平移集}得到.
以下设 $U\neq\mathbb{R}^n$.

\noindent$\mathbf{Step}\bm{1}$.{\heiti 确定 $\overline{V}$ 到边界的统一正距离.}

定义
\begin{align*}
d(x)\triangleq\operatorname{dist}(x,\partial U),
\quad x\in\mathbb{R}^n.
\end{align*}
由\cref{Real Analysis-theorem:距离函数都是一致连续的}知$d$ 连续.
又因为 $U$ 是开集, $\partial U\cap U=\varnothing$,由\subcref{Real Analysis-proposition:点集间的距离的性质}{Real Analysis-proposition:点集间的距离的性质-3},对每个 $x\in U$ 都有 $d(x)>0$.
因此, $d$ 在紧集 $\overline{V}\subseteq U$ 上取得正的最小值.

由于 $V$ 在 $\overline{V}$ 中稠密, 由 $d$ 的连续性有
\begin{align*}
\delta
=\inf\limits_{x\in V}d(x)
=\min\limits_{x\in\overline{V}}d(x)
>0.
\end{align*}
具体而言, 若 $x_0\in\overline{V}$ 是最小值点, 取 $x_k\in V$ 使 $x_k\to x_0$, 则
\begin{align*}
\inf\limits_{x\in V}d(x)
\leqslant
\lim\limits_{k\to\infty}d(x_k)
=d(x_0)
=\min\limits_{x\in\overline{V}}d(x),
\end{align*}
而反向不等式由$V\subseteq\overline{V}$ 立即得到.
特别地,
\begin{align*}
|x-z|\geqslant\delta,
\quad \forall x\in\overline{V},\, z\in\partial U.
\end{align*}

\noindent$\mathbf{Step}\bm{2}$.{\heiti 证明移动距离小于 $\delta$ 时不会离开 $U$.}

任取 $x\in\overline{V}$ 及满足 $|y-x|<\delta$ 的点 $y\in\mathbb{R}^n$.
下面证明 $y\in U$.
反设 $y\notin U$, 定义
\begin{align*}
\gamma(s)\triangleq x+s(y-x),
\quad s\in[0,1],
\end{align*}
以及
\begin{align*}
A\triangleq\{s\in[0,1]:\gamma(s)\notin U\}.
\end{align*}
由于 $\gamma$ 连续且 $\mathbb{R}^n\setminus U$ 是闭集,再由\cref{Topology-proposition:连续映射下闭集的原像是闭集}知$A=\gamma^{-1}(\mathbb{R}^n\setminus U)$ 是 $[0,1]$ 的闭子集.
又因为 $1\in A$, 所以 $A$ 是非空紧集, 从而存在最小值
\begin{align*}
s_0\triangleq\min A.
\end{align*}
由 $\gamma(0)=x\in U$, 有 $0\notin A$, 因此 $0<s_0\leqslant1$.
根据 $s_0$ 的最小性, 对所有 $0\leqslant s<s_0$ 都有 $\gamma(s)\in U$, 而 $\gamma(s_0)\notin U$.
由$\gamma$的连续性和\subcref{Real Analysis-theorem:边界和开集基本性质}{Real Analysis-theorem:边界和开集基本性质-5}知
\begin{align*}
z\triangleq\gamma(s_0)
=\lim\limits_{s\to s_0}\gamma(s)
\in\overline{U}\setminus U
=\partial U.
\end{align*}
于是
\begin{align*}
\delta
\leqslant|x-z|
=s_0|y-x|
\leqslant|y-x|
<\delta,
\end{align*}
矛盾!
因此 $y\in U$.

\noindent$\mathbf{Step}\bm{3}$.{\heiti 应用于线段平移, 并验证紧包含关系.}

若 $|a|<\delta$, 则对任意 $x\in\overline{V}$ 及 $t\in[0,1]$,
\begin{align*}
|(x+ta)-x|=t|a|\leqslant|a|<\delta.
\end{align*}
由$\mathbf{Step}\bm{2}$可知 $x+ta\in U$, 因而
\begin{align*}
K_a\triangleq
\{x+ta:x\in\overline{V},\, t\in[0,1]\}
\subseteq U.
\end{align*}
由\cref{proposition:紧集平移与紧参数平移集}, $K_a$ 是紧集.
所以 $K_a\subset\subset U$.
同理, $\overline{V}+a$ 是包含于 $U$ 的紧集, 故 $\overline{V}+a\subset\subset U$.

最后, 对有限个满足 $|a_j|<\delta$ 的向量分别应用上述结论, 可知每个 $K_{a_j}$ 都是包含于 $U$ 的紧集.
它们的有限并仍然是紧集且包含于 $U$, 因此
\begin{align*}
\bigcup\limits_{j=1}^m K_{a_j}
=
\bigcup\limits_{j=1}^m
\{x+ta_j:x\in\overline{V},\, t\in[0,1]\}
\subset\subset U.
\end{align*}
\end{proof}

\begin{theorem}[差商的分部积分公式]\label{theorem:ch05-差商的分部积分公式}
设 $U\subseteq\mathbb{R}^n$ 为开集, $i\in\{1,\cdots,n\}$, $h\neq 0$. 设 $u\in L^1_{\mathrm{loc}}(U)$, $v\in C_c^\infty(U)$, 且
\begin{align*}
0<|h|<\frac{1}{2}\operatorname{dist}(\operatorname{spt}v,\partial U),
\end{align*}
则
\begin{align*}
\int_U u(x)D_i^h v(x)\,\mathrm{d}x
=-\int_U D_i^{-h}u(x)v(x)\,\mathrm{d}x.
\end{align*}
\end{theorem}
\begin{proof}
由 $v\in C_c^\infty(U)$ 及 $0<|h|<\frac{1}{2}\operatorname{dist}(\operatorname{spt}v,\partial U)$,再结合\cref{proposition:开集内部小幅平移后仍在原本开集中}可知对任意 $x\in\operatorname{spt}v$, 有 $x+h e_i\in U$ 且 $x-h e_i\in U$. 因此下列积分均有意义. 直接计算得
\begin{align*}
\int_U u(x)D_i^h v(x)\,\mathrm{d}x
=\frac{1}{h}\int_U u(x)\left(v(x+h e_i)-v(x)\right)\mathrm{d}x=\frac{1}{h}\int_U u(x)v(x+h e_i)\,\mathrm{d}x
-\frac{1}{h}\int_U u(x)v(x)\,\mathrm{d}x.
\end{align*}
对第一个积分作变量替换 $y=x+h e_i$, 则 $x=y-h e_i$, $\mathrm{d}x=\mathrm{d}y$. 由于 $v$ 的支集在 $U$ 中, 当 $y\notin U$ 时 $v(y)=0$, 故积分区域仍可取为 $U$, 从而
\begin{align*}
\int_U u(x)v(x+h e_i)\,\mathrm{d}x
=\int_U u(y-h e_i)v(y)\,\mathrm{d}y.
\end{align*}
代回上式, 得
\begin{align*}
\int_U u(x)D_i^h v(x)\,\mathrm{d}x
&=\frac{1}{h}\int_U u(x-h e_i)v(x)\,\mathrm{d}x
-\frac{1}{h}\int_U u(x)v(x)\,\mathrm{d}x\\
&=-\int_U \frac{u(x)-u(x-h e_i)}{h}v(x)\,\mathrm{d}x\\
&=-\int_U D_i^{-h}u(x)v(x)\,\mathrm{d}x.
\end{align*}
\end{proof}

\begin{theorem}[差商与弱导数]\label{theorem:ch05-差商弱导数}
设$U$是$\mathbb{R}^n$中的开集,对每个 $V\subset \subset U$.
\begin{enumerate}[(i)]
\item 设 $1\leqslant p<\infty$，$u\in W^{1,p}(U)$。则对一切 $0<|h|<\frac{1}{2}\operatorname{dist}(V,\partial U)$,都有
\begin{align}
\|D^{h}u\|_{L^{p}(V)}\leqslant n\|Du\|_{L^{p}(U)}
\end{align}

\item 设 $1<p<\infty$，$u\in L^{p}(V)$，且存在常数 $C$ 使
\begin{align}\label{eq::jijiowjiwiw23323rwe}
\|D^{h}u\|_{L^{p}(V)}\leqslant C
\end{align}
对一切 $0<|h|<\frac{1}{2}\operatorname{dist}(V,\partial U)$ 成立。则
\begin{align*}
u\in W^{1,p}(V),\quad
\|Du\|_{L^p(V)}\leqslant C.
\end{align*}
\end{enumerate}
\end{theorem}
\begin{remark}
$p=\infty$的情形见下文\cref{theorem:ch05-W1infty刻画}.断言(ii)对 $p=1$ 不成立。
下面的一维反例说明, 差商在 $L^1(V)$ 中一致有界, 并不足以推出 $u\in W^{1,1}(V)$.

取
\begin{align*}
U=(-2,2),\quad V=(-1,1),
\end{align*}
并定义
\begin{align*}
u(x)\triangleq
\begin{cases}
0,&-2<x\leqslant0,\\
1,&0<x<2.
\end{cases}
\end{align*}
显然 $V\subset \subset U$, 且 $u\in L^1(U)$. 又因为
\begin{align*}
\partial U=\{-2,2\},\quad
\operatorname{dist}(V,\partial U)=1,
\end{align*}
所以题设中的步长条件就是
$0<|h|<\frac{1}{2}.$

先计算正步长的差商. 当 $0<h<\frac{1}{2}$ 时, 只有取值点 $x$ 和 $x+h$ 分别位于跳跃点 $0$ 的两侧, 差商才可能非零. 具体而言, 对几乎处处的 $x\in V$, 有
\begin{align*}
D^h u(x)
=
\frac{u(x+h)-u(x)}{h}
=
\begin{cases}
\dfrac{1}{h},&-h<x<0,\\
0,&x\in V\setminus(-h,0).
\end{cases}
\end{align*}
因此
\begin{align*}
\|D^h u\|_{L^1(V)}
=
\int_{-h}^{0}\frac{1}{h}\,\mathrm{d}x
=1.
\end{align*}

再计算负步长的差商. 当 $-\frac{1}{2}<h<0$ 时,对几乎处处的 $x\in V$, 有
\begin{align*}
D^h u(x)
=
\frac{u(x+h)-u(x)}{h}
=
\frac{u(x)-u(x-(-h))}{-h}=
\begin{cases}
-\frac{1}{h},&0<x<-h,\\
0,&x\in V\setminus(0,-h).
\end{cases}
\end{align*}
从而
\begin{align*}
\|D^h u\|_{L^1(V)}
=
\int_0^{-h}-\frac{1}{h}\,\mathrm{d}x
=1.
\end{align*}
综合正、负步长的计算, 对一切
$0<|h|<\frac{1}{2}\operatorname{dist}(V,\partial U)$, 都有
\begin{align*}
\|D^h u\|_{L^1(V)}=1.
\end{align*}
因此, 题设中的一致有界条件成立, 可以取 $C=1$.

下面证明 $u\notin W^{1,1}(V)$.
反设 $u\in W^{1,1}(V)$, 则存在弱导数 $g\in L^1(V)$, 使得对任意 $\varphi\in C_c^\infty(V)$, 都有
\begin{align*}
\int_V u(x)\varphi'(x)\,\mathrm{d}x
=
-\int_V g(x)\varphi(x)\,\mathrm{d}x.
\end{align*}
另一方面, 由$\text{spt}\,\varphi\subseteq V=(-1,1)$知$\varphi$ 在端点 $1$ 附近恒为零,再根据 $u$ 的定义,有
\begin{align*}
\int_V u(x)\varphi'(x)\,\mathrm{d}x
=
\int_0^1\varphi'(x)\,\mathrm{d}x
=
-\varphi(0).
\end{align*}
因此, 假定存在的弱导数 $g$ 必须满足
\begin{align}\label{NMZNINWhiuhfjknxlzzvz}
\int_V g(x)\varphi(x)\,\mathrm{d}x
=
\varphi(0),
\quad
\forall\varphi\in C_c^\infty(V).
\end{align}

取标准磨光子 $\eta\in C_c^\infty(\mathbb{R}^n)$, 满足
\begin{align*}
\eta\geqslant0,\qquad
\operatorname{spt}\eta\subseteq\overline{B^{0}(0,1)},\qquad
\int_{\mathbb{R}^n}\eta(x)\,\mathrm{d}x=1.
\end{align*}
对$\varepsilon>0$, 定义
\begin{align*}
\eta_\varepsilon(x)
\triangleq\frac{1}{\varepsilon^n}
\eta\left(\frac{x}{\varepsilon}\right).
\end{align*}
于是
\begin{align*}
\eta(0)>0,\quad
0\leqslant\eta(x)\leqslant\eta(0),
\quad x\in\mathbb{R}.
\end{align*}
对 $0<r<1$, 直接取测试函数
\begin{align*}
\varphi_r(x)
\triangleq
\frac{\eta_{\frac{r}{2}}(x)}{\eta_{\frac{r}{2}}(0)}
=
\frac{\eta\left(\frac{2x}{r}\right)}{\eta(0)}.
\end{align*}
由于 $\eta\in C_c^\infty(\mathbb{R})$, 且
$\operatorname{spt}\eta=[-1,1]$, 有
\begin{align*}
\varphi_r\in C^\infty(\mathbb{R}),\quad
\operatorname{spt}\varphi_r
=
\left[-\frac{r}{2},\frac{r}{2}\right]
\subset \subset (-r,r)\subset \subset V.
\end{align*}
因此
\begin{align*}
\varphi_r\in C_c^\infty(V),\quad
0\leqslant\varphi_r\leqslant1,\quad
\varphi_r(0)=1.
\end{align*}
将 $\varphi_r$ 代入\eqref{NMZNINWhiuhfjknxlzzvz}式, 得到
\begin{align}\label{eqNVMZnnjczvnms}
1
=
\left|
\int_V g(x)\varphi_r(x)\,\mathrm{d}x
\right|
\leqslant
\int_{-\frac{r}{2}}^{\frac{r}{2}}|g(x)|\,\mathrm{d}x
\leqslant
\int_{-r}^{r}|g(x)|\,\mathrm{d}x.
\end{align}
但是, 由于 $g\in L^1(V)$, 由\hyperref[Real Analysis-theorem:积分的绝对连续性]{Lebesgue积分的绝对连续性}, 有
\begin{align*}
\lim\limits_{r\to0^+}
\int_{-r}^{r}|g(x)|\,\mathrm{d}x
=0,
\end{align*}
这与上面的不等式\eqref{eqNVMZnnjczvnms}矛盾! 因此不存在这样的 $g\in L^1(V)$, 即
\begin{align*}
u\notin W^{1,1}(V).
\end{align*}

这个反例中, 差商的非零部分位于长度为 $|h|$ 的小区间上, 高度为 $\frac{1}{|h|}$, 因而其 $L^1$ 范数始终等于 $1$.
但是当 $h\to0$ 时, 这些差商的积分质量越来越集中于跳跃点 $0$.
这正是本例中差商虽然在 $L^1(V)$ 中一致有界, 却不能产生 $L^1$ 弱导数的原因.

直观地看, 每个差商 $D^h u$ 都是一个紧贴跳跃点 $0$ 的矩形脉冲: 它的宽度为 $|h|$, 高度为 $\frac{1}{|h|}$, 所以矩形的面积始终为 $1$, 即
\begin{align*}
\int_V|D^h u(x)|\,\mathrm{d}x
=|h|\cdot\frac{1}{|h|}
=1.
\end{align*}
当 $h\to0$ 时, 这个矩形越来越窄, 同时越来越高, 但面积并没有减少. 因此, 对每个固定的 $x\neq0$, 当 $|h|$ 足够小时都有 $D^h u(x)=0$, 可是差商的总积分仍然等于 $1$. 这说明积分质量并没有消失, 而是被压缩到越来越靠近 $0$ 的小区间中.

所谓积分质量集中于 $0$, 具体是指: 无论取多小的邻域 $(-r,r)$, 只要 $0<|h|<\min\{r,\frac{1}{2}\}$, 差商的全部积分质量就已经落在这个邻域内, 即
\begin{align*}
\int_{-r}^{r}|D^h u(x)|\,\mathrm{d}x=1.
\end{align*}
如果$u$存在弱导数$g$,那么$g$的积分性质也应该与$D^hu(x)$类似.
而任何固定的 $L^1$ 函数 $g$ 都不可能具有这种性质. 因为 Lebesgue 积分的绝对连续性保证
\begin{align*}
\lim\limits_{r\to 0^+}\int_{-r}^{r}|g(x)|\,\mathrm{d}x=0.
\end{align*}
也就是说, 对于一个固定的可积函数, 区间缩得足够小时, 其中的积分质量必然趋于零; 本例中的差商却通过不断增大高度, 在不断缩小的区间内始终保留总量为 $1$ 的积分质量.

函数 $u$ 在 $0$ 两侧都是常数, 它的变化全部来自 $0$ 处大小为 $1$ 的跳跃. 前面的测试函数计算说明, 若要用一个 $L^1$ 弱导数表示这个跳跃, 就必须让它在任意小的零点邻域内保留至少为 $1$ 的绝对积分, 这与积分的绝对连续性矛盾！ 因此, 差商的 $L^1$ 范数一致有界只能控制积分质量的总量, 不能阻止这些质量集中到一个点上, 也就不足以保证存在 $L^1$ 弱导数.
\end{remark}
\begin{proof}
\noindent$\mathbf{Step}\bm{1}$.{\heiti 选择包含所有平移线段的中间区域, 并建立光滑函数的局部差商估计.}

先证明断言(i).固定任意一个满足 $0<|h|<\frac{\operatorname{dist}(V,\partial U)}{2}$ 的步长 $h$, 定义
\begin{align*}
K_h\triangleq
\bigcup\limits_{i=1}^n
\left\{
x+the_i:
x\in\overline{V},\, t\in[0,1]
\right\}.
\end{align*}
由\cref{proposition:紧集平移与紧参数平移集}知$K_h$也是紧集.
又由\cref{proposition:开集内部小幅平移后仍在原本开集中}知$K_h\subseteq U$.

由\cref{proposition:紧集的相对紧开邻域}, 可以选择开集 $W$, 使得
\begin{align*}
K_h\subseteq W\subset\subset U.
\end{align*}
特别地,有
\begin{align}\label{eqNJKNDnxjklndklszjsjwzzcxvds}
V+the_i\subseteq W,
\quad t\in[0,1],\quad i=1,\cdots,n.
\end{align}

设 $u\in C^\infty(W)\cap W^{1,p}(W)$.
对任意 $x\in V$ 及 $i=1,\cdots,n$, 线段 $\{x+the_i:t\in[0,1]\}$ 完全位于 $W$ 中, 因而由微积分基本定理有
\begin{align*}
u\left( x+he_i \right) -u\left( x \right) =\int_0^1{\frac{\mathrm{d}}{\mathrm{d}t}u\left( x+the_i \right) \,\mathrm{d}t}=h\int_0^1{u_{x_i}\left( x+the_i \right) \,\mathrm{d}t}\Longrightarrow D^h_i u(x)=\int_0^1{u_{x_i}\left( x+the_i \right) \,\mathrm{d}t}.
\end{align*}
从而再由\hyperref[Basis of Analytics-theorem:离散的Hölder不等式]{离散的Hölder不等式}得
\begin{align*}
|D^h u(x)|^p
&\leqslant \left( \sum_{i=1}^n |D_i^h u(x)| \right)^p
\leqslant n^{p-1} \sum_{i=1}^n |D_i^h u(x)|^p \\
&\leqslant n^{p-1} \sum_{i=1}^n \int_0^1 |u_{x_i}(x + t h e_i)|^p \,\mathrm{d}t \\
&\leqslant n^{p-1} \sum_{i=1}^n \int_0^1 |Du(x + t h e_i)|^p \,\mathrm{d}t.
\end{align*}
两边在 $V$ 上积分并利用\hyperref[Real Analysis-theorem:实变函数-Fubini定理]{Fubini定理}和\eqref{eqNJKNDnxjklndklszjsjwzzcxvds}式,即得
\begin{align*}
\int_V{|D^hu|^p\,\mathrm{d}x}&\leqslant n^{p-1}\sum_{i=1}^n{\int_0^1{\int_V{\left| Du\left( x+the_i \right) \right|^p\,\mathrm{d}x\,\mathrm{d}t}}}=n^{p-1}\sum_{i=1}^n{\int_0^1{\int_{V+the_i}{\left| Du\left( x \right) \right|^p\,\mathrm{d}x\,\mathrm{d}t}}}
\\
&\leqslant n^{p-1}\sum_{i=1}^n{\int_0^1{\int_W{\left| Du\left( x \right) \right|^p\,\mathrm{d}x\,\mathrm{d}t}}}=n^p\int_W{\left| Du\left( x \right) \right|^p\,\mathrm{d}x}.
\end{align*}
故
\begin{align}\label{mzvjkhshtuwwaa}
\|D^h u_m\|_{L^p(V)}\leqslant n\|Du_m\|_{L^p(W)}.
\end{align}

\noindent$\mathbf{Step}\bm{2}$.{\heiti 利用局部光滑逼近将估计推广到 Sobolev 函数.}

仍固定上述 $h$ 和相应的 $W$.
由于 $W\subset\subset U$, 由\hyperref[theorem:ch05-光滑函数的局部逼近定理]{光滑函数的局部逼近定理}, 可以选择一列磨光函数 $u_m$, 使得
\begin{align*}
u_m\in C^\infty(W)\cap W^{1,p}(W),\quad
\lim\limits_{m\to\infty}
\|u_m-u\|_{W^{1,p}(W)}=0.
\end{align*}
于是 $u_m\to u$ 于 $L^p(W)$，且 $Du_m\to Du$ 于 $L^p(W)$。
对每个 $i=1,\cdots,n$, 由 $V\subseteq W$ 及 $V+he_i\subseteq W$, 有
\begin{align*}
\|D_i^h u_m-D_i^h u\|_{L^p(V)}
&\leqslant
\frac{1}{|h|}
\left(
\|(u_m-u)(\,\cdot+he_i)\|_{L^p(V)}
+\|u_m-u\|_{L^p(V)}
\right)\\
&=
\frac{1}{|h|}
\left(
\|u_m-u\|_{L^p(V+he_i)}
+\|u_m-u\|_{L^p(V)}
\right)\\
&\leqslant
\frac{2}{|h|}\|u_m-u\|_{L^p(W)}.
\end{align*}
因此, 对整个向量差商,
\begin{align*}
\|D^h u_m-D^h u\|_{L^p(V)}
\leqslant
\sum\limits_{i=1}^n
\|D_i^h u_m-D_i^h u\|_{L^p(V)}
\leqslant
\frac{2n}{|h|}\|u_m-u\|_{L^p(W)}
\longrightarrow0,
\quad m\to\infty.
\end{align*}
故 $D^h u_m\to D^h u$ 于 $L^p(V)$。
由于 $u_m$ 光滑，利用$\mathbf{Step}\bm{1}$中的\eqref{mzvjkhshtuwwaa}式可得
\begin{align*}
\|D^h u_m\|_{L^p(V)}\leqslant n\|Du_m\|_{L^p(W)}.
\end{align*}
利用范数的连续性, 令上式$m\to\infty$, 可得
\begin{align*}
\|D^h u\|_{L^p(V)}
\leqslant n\|Du\|_{L^p(W)}\leqslant n\|Du\|_{L^p(U)}.
\end{align*}
因而断言(i)成立.

\noindent$\mathbf{Step}\bm{3}$.{\heiti 利用一致有界性, 对向量差商提取弱收敛子列.}

下面证明断言(ii).假设\eqref{eq::jijiowjiwiw23323rwe}蕴含
\begin{align*}
\sup_{0<h<\frac{\operatorname{dist}(V,\partial U)}{2}}\|D_{i}^{-h}u\|_{L^{p}(V)}<\infty;
\end{align*}
又由$1<p<\infty$和\cref{Functional Analysis-proposition:Lp对偶与自反性}知$L^p(V)$是自反空间，再根据\cref{Functional Analysis-proposition:自反空间中弱列紧性等价与有界},存在函数 $v_{i}\in L^{p}(V)$ 及子列 $h_{k}\to 0$ 使
\begin{align}\label{eqr::JIjiwjijjJFEW8334435JIOFJA-2}
D_{i}^{-h_{k}}u\rightharpoonup v_{i}\quad\text{弱收敛于}\ L^{p}(V).
\end{align}

\noindent$\mathbf{Step}\bm{4}$.{\heiti 利用差商分部积分识别弱导数并得到范数估计.}

固定 $i\in\{1,\cdots,n\}$, 任取 $\phi\in C_c^\infty(V)$, 并将 $\phi$ 在 $V$ 外延拓为零.
由于 $\phi$ 的支集紧包含于 $V$, 对充分小的 $h$, 由\hyperref[theorem:ch05-差商的分部积分公式]{差商的分部积分公式}, 有
\begin{align}
\int_VuD_i^h\phi\,\mathrm{d}x
=
-\int_V(D_i^{-h}u)\phi\,\mathrm{d}x.
\label{eqr::JIjiwjijjJFEW8334435JIOFJA-1}
\end{align}
由于 $\phi$ 光滑且具有紧支集,再结合偏导数定义知$D_i^{h_k}\phi\to\phi_{x_i}$ 一致收敛, 从而再结合\eqref{eqr::JIjiwjijjJFEW8334435JIOFJA-2}\eqref{eqr::JIjiwjijjJFEW8334435JIOFJA-1}式,得到
\begin{align*}
\int_Vu\phi_{x_i}\,\mathrm{d}x
=
\lim\limits_{k\to\infty}
\int_VuD_i^{h_k}\phi\,\mathrm{d}x
=
-\lim\limits_{k\to\infty}
\int_V(D_i^{-h_k}u)\phi\,\mathrm{d}x
=
-\int_Vv_i\phi\,\mathrm{d}x.
\end{align*}
因此 $v_i=u_{x_i}$ 是 $u$ 在 $V$ 中的第 $i$ 个弱偏导数（$i=1,\cdots,n$）.
由于每个 $v_i\in L^p(V)$, 且 $u\in L^p(V)$, 可得
\begin{align*}
u\in W^{1,p}(V),\quad Du=(v_1,\cdots,v_n).
\end{align*}
由\hyperref[Functional Analysis-proposition:范数关于弱收敛的下半连续性]{范数关于弱收敛的下半连续性}得
\begin{align*}
\|Du\|_{L^p(V)}
\leqslant
\varliminf\limits_{k\to\infty}
\|D^{-h_k}u\|_{L^p(V)}
\leqslant C.
\end{align*}
故断言(ii)成立.
\end{proof}

\begin{proposition}[平直边界附近的切向差商估计]\label{proposition:ch05-平直边界附近的切向差商估计}
设 $n\geqslant2$, $0<r<R<\infty$, 并记
\begin{align*}
U\triangleq B^{0}(0,R)\cap\{x_n>0\},\quad
V\triangleq B^{0}(0,r)\cap\{x_n>0\}.
\end{align*}
设 $1\leqslant p<\infty$, $u\in W^{1,p}(U)$. 对每个切向方向 $i=1,\cdots,n-1$ 以及每个满足 $0<|h|<R-r$ 的步长 $h$, 差商
\begin{align*}
D_i^h u(x)\triangleq\frac{u(x+he_i)-u(x)}{h}
\end{align*}
在 $V$ 上几乎处处有定义, 且满足
\begin{align}\label{eq:ch05-TanDQ-r7K2m9P4}
\int_V|D_i^h u(x)|^p\,\mathrm{d}x
\leqslant
\int_U|u_{x_i}(x)|^p\,\mathrm{d}x.
\end{align}
这里 $e_i$ 是第 $i$ 个标准基向量.

特别地, 当 $R=1$, $r=\frac{1}{2}$ 时, 尽管 $V\subset\subset U$ 不成立, 估计 \eqref{eq:ch05-TanDQ-r7K2m9P4} 对所有 $i=1,\cdots,n-1$ 和 $0<|h|<\frac{1}{2}$ 仍然成立.
\end{proposition}

\begin{proof}
固定 $i\in\{1,\cdots,n-1\}$ 以及满足 $0<|h|<R-r$ 的步长 $h$.

\noindent$\mathbf{Step}\bm{1}$.{\heiti 确认切向平移后的线段始终位于 $U$ 内.}

对任意 $x\in V$ 和 $t\in[0,1]$, 由于 $i\leqslant n-1$, 切向平移不改变第 $n$ 个坐标, 因而
\begin{align*}
(x+the_i)_n=x_n>0.
\end{align*}
另一方面,
\begin{align*}
|x+the_i|
\leqslant |x|+t|h|
<r+|h|
<R.
\end{align*}
所以
\begin{align*}
x+the_i\in U,\quad
x\in V,\quad t\in[0,1].
\end{align*}
特别地, 对所有 $t\in[0,1]$, 都有
\begin{align*}
V+the_i\subseteq U.
\end{align*}
因此, 无论 $h$ 为正还是为负, 差商中涉及的取值点都位于 $U$ 内.

\noindent$\mathbf{Step}\bm{2}$.{\heiti 估计光滑函数的差商.}

先设 $u\in C^\infty(U)\cap W^{1,p}(U)$. 对任意 $x\in V$, 由微分学基本定理得
\begin{align*}
u(x+he_i)-u(x)
=\int_0^1\frac{\mathrm{d}}{\mathrm{d}t}u(x+the_i)\,\mathrm{d}t
=h\int_0^1u_{x_i}(x+the_i)\,\mathrm{d}t.
\end{align*}
除以 $h$, 得
\begin{align*}
D_i^h u(x)
=\int_0^1u_{x_i}(x+the_i)\,\mathrm{d}t.
\end{align*}
当 $p>1$ 时, 对长度为 $1$ 的区间 $[0,1]$ 使用\hyperref[Real Analysis-theorem:实变函数--Hölder不等式]{Hölder不等式}; 当 $p=1$ 时, 直接使用积分的绝对值不等式. 两种情形均给出
\begin{align*}
|D_i^h u(x)|^p
\leqslant
\left(\int_0^1|u_{x_i}(x+the_i)|\,\mathrm{d}t\right)^p
\leqslant
\int_0^1|u_{x_i}(x+the_i)|^p\,\mathrm{d}t.
\end{align*}
两边在 $V$ 上积分, 利用\hyperref[Real Analysis-theorem:实变函数-Fubini定理]{Fubini定理}和平移换元, 得
\begin{align}
\int_V|D_i^h u(x)|^p\,\mathrm{d}x
&\leqslant
\int_0^1\int_V|u_{x_i}(x+the_i)|^p\,\mathrm{d}x\,\mathrm{d}t=
\int_0^1\int_{V+the_i}|u_{x_i}(y)|^p\,\mathrm{d}y\,\mathrm{d}t\nonumber\\
&\leqslant
\int_0^1\int_U|u_{x_i}(y)|^p\,\mathrm{d}y\,\mathrm{d}t
=
\int_U|u_{x_i}(y)|^p\,\mathrm{d}y.\label{eq::NKZNJIOIJaaaaa}
\end{align}
这就证明了光滑情形下的结论.

\noindent$\mathbf{Step}\bm{3}$.{\heiti 截去边界薄层,利用局部光滑逼近.}

现在设 $u\in W^{1,p}(U)$. 为使用局部光滑逼近, 对每个 $0<\varepsilon<r$, 定义
\begin{align*}
V_\varepsilon
\triangleq V\cap\{x_n>\varepsilon\}.
\end{align*}
考虑从 $\overline{V_\varepsilon}$ 出发的所有平移线段组成的集合
\begin{align*}
K_\varepsilon
\triangleq
\left\{
x+the_i:
x\in\overline{V_\varepsilon},\, t\in[0,1]
\right\}.
\end{align*}
由\cref{proposition:紧集平移与紧参数平移集}知$K_\varepsilon$ 是紧集. 对任意 $y=x+the_i\in K_\varepsilon$, 有
\begin{align*}
y_n=x_n\geqslant\varepsilon>0,\quad
|y|\leqslant r+|h|<R.
\end{align*}
因此 $K_\varepsilon$ 是 $U$ 的紧子集. 由\cref{proposition:紧集的相对紧开邻域},可以选取开集 $W_\varepsilon$, 使得
\begin{align*}
K_\varepsilon\subseteq W_\varepsilon\subset\subset U.
\end{align*}
于是
\begin{align*}
V_\varepsilon+the_i\subseteq W_\varepsilon,\quad t\in[0,1].
\end{align*}

由\hyperref[theorem:ch05-光滑函数的局部逼近定理]{光滑函数的局部逼近定理}, 存在
$u_m\in C^\infty(W_\varepsilon)\cap W^{1,p}(W_\varepsilon)$, 使得
\begin{align*}
\|u_m-u\|_{W^{1,p}(W_\varepsilon)}
\longrightarrow0,\quad m\to\infty.
\end{align*}
由于从 $V_\varepsilon$ 出发的上述平移线段全部包含于 $W_\varepsilon$, 将\eqref{eq::NKZNJIOIJaaaaa}式的计算应用于 $u_m$, 得
\begin{align}\label{NJKHiwiiwwww}
\int_{V_\varepsilon}|D_i^h u_m(x)|^p\,\mathrm{d}x
\leqslant
\int_0^1\int_{V_\varepsilon+the_i}
|(u_m)_{x_i}(y)|^p\,\mathrm{d}y\,\mathrm{d}t
\leqslant
\int_{W_\varepsilon}|(u_m)_{x_i}(y)|^p\,\mathrm{d}y.
\end{align}

下面说明可以在此估计中令 $m\to\infty$. 因为 $h$ 已经固定, 且
$V_\varepsilon,V_\varepsilon+he_i\subseteq W_\varepsilon$, 所以
\begin{align*}
\|D_i^h u_m-D_i^h u\|_{L^p(V_\varepsilon)}
&\leqslant
\frac{1}{|h|}
\left(
\|u_m-u\|_{L^p(V_\varepsilon+he_i)}
+
\|u_m-u\|_{L^p(V_\varepsilon)}
\right)\\
&\leqslant
\frac{2}{|h|}
\|u_m-u\|_{L^p(W_\varepsilon)}
\longrightarrow0,\quad m\to\infty.
\end{align*}
同时,还有
\begin{align*}
\|(u_m)_{x_i}-u_{x_i}\|_{L^p(W_\varepsilon)}\leqslant \|u_m-u\|_{L^p(W_\varepsilon)}
\longrightarrow0,\quad m\to\infty.
\end{align*}
因此, 由 $L^p$ 范数的连续性, 在\eqref{NJKHiwiiwwww}式中令 $m\to\infty$, 得
\begin{align*}
\int_{V_\varepsilon}|D_i^h u(x)|^p\,\mathrm{d}x
\leqslant
\int_{W_\varepsilon}|u_{x_i}(x)|^p\,\mathrm{d}x\leqslant
\int_U|u_{x_i}(x)|^p\,\mathrm{d}x.
\end{align*}

\noindent$\mathbf{Step}\bm{4}$.{\heiti 令边界薄层的厚度趋于零.}

当 $\varepsilon\to0^+$ 时, 集合 $V_\varepsilon$ 单调增大趋于$V$. 由\hyperref[Real Analysis-theorem:Beppo Levi非负渐升列积分定理]{Levi定理},得
\begin{align*}
\int_V{|D_{i}^{h}u(x)|^p\,\mathrm{d}x}&=\int_{\mathbb{R} ^n}{|D_{i}^{h}u(x)|^p\chi _V\left( x \right) \,\mathrm{d}x}=\lim_{\varepsilon \rightarrow 0^+} \int_{\mathbb{R} ^n}{|D_{i}^{h}u(x)|^p\chi _{V_{\varepsilon}}\left( x \right) \,\mathrm{d}x}
\\
&=\lim_{\varepsilon \rightarrow 0^+} \int_{V_{\varepsilon}}{|D_{i}^{h}u(x)|^p\,\mathrm{d}x}\leqslant \int_U{|u_{x_i}(x)|^p\,\mathrm{d}x}.
\end{align*}
由于方向 $i\in\{1,\cdots,n-1\}$ 和步长 $0<|h|<R-r$ 是任意选取的, \eqref{eq:ch05-TanDQ-r7K2m9P4}式对所有这些方向和步长均成立.
\end{proof}


\begin{theorem}[$W^{1,\infty}_{\mathrm{loc}}$ 的刻画]\label{theorem:ch05-W1infty局部刻画}
设 $U \subseteq \mathbb{R}^n$ 为开集，$f: U \to \mathbb{R}$。则
$f\ \text{在}\ U\ \text{中局部 Lipschitz 连续当且仅当}f \in W^{1,\infty}_{\mathrm{loc}}(U).$
特别地,$f$的Lipschitz 常数不超过 \(\|Df\|_{L^\infty(W)}\)
\end{theorem}
\begin{proof}

{\heiti 必要性:}
设 $f$ 在 $U$ 中局部 Lipschitz 连续 . 固定 $i\in\{1,\cdots,n\}$ . 任取开集 $V \subset\subset W \subset\subset U$ , 取 $0<|h|<\operatorname{dist}(V,\partial W)$ , 定义差商
\begin{align}\label{eq:局部差商定义}
D_{i}^{-h}f(x)\triangleq\frac{f(x)-f(x-he_i)}{h},\qquad x\in V.
\end{align}
由 $f$ 在 $W$ 上的 Lipschitz 条件 , 有
\begin{align*}
\sup_{0<|h|<\operatorname{dist}(V,\partial W)}\|D_{i}^{-h}f\|_{L^\infty(V)}
\leqslant \operatorname{Lip}(f|_W)<\infty.
\end{align*}
因此差商族 $\{D_{i}^{-h}f\}_{h\neq 0}$ 在 $L^\infty(V)$ 中一致有界 , 从而在 $L^2(V)$ 中也一致有界 . 由\cref{Functional Analysis-proposition:Lp对偶与自反性}知 $L^2(V)$ 是自反的 , 再根据\cref{Functional Analysis-proposition:自反空间中弱列紧性等价与有界} , 存在子列 $h_k\to 0$ 及函数 $g_i\in L^2(V)$ , 使得
\begin{align*}
D_{i}^{-h_k}f\rightharpoonup g_i\quad\text{在}\ L^2(V)\ \text{中弱收敛}.
\end{align*}
另一方面，由 $\|D_i^{-h}f\|_{L^\infty(\mathbb{R}^n)}\leqslant \operatorname{Lip}(f|_W)$ 对一切 $h\neq 0$ 成立，得
\begin{align*}
\left|\int_{\mathbb{R}^n}\phi D_i^{-h_k}f\,\mathrm{d}x\right|
\leqslant \|D_i^{-h_k}f\|_{L^\infty(\mathbb{R}^n)}\int_{\mathbb{R}^n}|\phi|\,\mathrm{d}x
\leqslant \operatorname{Lip}(f|_W)\int_{\mathbb{R}^n}|\phi|\,\mathrm{d}x.
\end{align*}
令 $k\to\infty$，即得
\begin{align*}
\left|\int_{\mathbb{R}^n}g_i\,\phi\,\mathrm{d}x\right|
\leqslant \operatorname{Lip}(u)\int_{\mathbb{R}^n}|\phi|\,\mathrm{d}x.
\end{align*}
由于\cref{Real Analysis-theorem:L无穷空间范数对偶刻画}以及$C_c^\infty(\mathbb{R}^n)$ 在 $L^1(\mathbb{R}^n)$ 中稠密，有
\begin{align*}
\|g_i\|_{L^\infty(\mathbb{R}^n)}
=\sup_{\substack{\phi\in C_c^\infty(\mathbb{R}^n)\\ \|\phi\|_{L^1(\mathbb{R}^n)}= 1}}
\left|\int_{\mathbb{R}^n}g_i\,\phi\,\mathrm{d}x\right|.
\end{align*}
因此综上可得
\begin{align}\label{eq:局部弱极限L无穷界}
\|g_i\|_{L^\infty(V)}\leqslant \operatorname{Lip}(f|_W).
\end{align}

任取 $\phi\in C_c^\infty(V)$ . 由\hyperref[theorem:ch05-差商的分部积分公式]{差商的分部积分公式} , 当$h$充分小时,有
\begin{align}\label{eq:局部差商分部积分}
\int_{V}f(x)D_{i}^{h}\phi(x)\,\mathrm{d}x
=-\int_{V}D_{i}^{-h}f(x)\phi(x)\,\mathrm{d}x.
\end{align}
其中 $D_{i}^{h}\phi(x)=\frac{\phi(x+he_i)-\phi(x)}{h}$ . 因为 $\phi$ 光滑且具有紧支集 , 所以当 $h_k\to 0$ 时 , $D_{i}^{h_k}\phi\to \phi_{x_i}$ 在 $L^\infty(V)$ 中成立 . 对 \eqref{eq:局部差商分部积分} 式令 $h=h_k\to 0$ , 左端由\hyperref[Real Analysis-theorem:控制收敛定理]{控制收敛定理}收敛到 $\int_V f\phi_{x_i}\,\mathrm{d}x$ ; 右端由弱收敛定义收敛到 $-\int_V g_i\phi\,\mathrm{d}x$ . 于是
\begin{align*}
\int_{V}f\phi_{x_i}\,\mathrm{d}x=-\int_{V}g_i\phi\,\mathrm{d}x.
\end{align*}
上式对一切 $\phi\in C_c^\infty(V)$ 成立 , 故 $g_i$ 是 $f$ 关于 $x_i$ 的弱导数 , 即 $g_i=f_{x_i}$ 在弱意义下成立 . 由于 $g_i\in L^\infty(V)$ , 且这对任意 $V\subset\subset U$ 均成立 , 所以 $f\in W^{1,\infty}_{\mathrm{loc}}(U)$ .

{\heiti 充分性：}
设 \(f\in W^{1,\infty}_{\mathrm{loc}}(U)\)。任取闭球 \(B\subset\subset U\)，取开集 \(W\) 使得 \(B\subset\subset W\subset\subset U\)。令 \(f^\varepsilon\triangleq\eta_\varepsilon*f\) 为通常的磨光函数，其中 \(\varepsilon<\operatorname{dist}(B,\partial W)\)。
由\hyperref[theorem:appC-光滑算子的性质]{光滑算子的性质}知 \(f^\varepsilon\in C^\infty(W)\)，且对任意 \(i=1,\cdots,n\)，有
\begin{align*}
D_if^\varepsilon(x)=(\eta_\varepsilon*D_if)(x).
\end{align*}
根据\hyperref[theorem:appC-光滑算子的性质]{光滑算子的性质}，有
\begin{align}\label{eq:磨光导数一致界}
\|Df^\varepsilon\|_{L^\infty(B)}\leqslant \|Df\|_{L^\infty(W)}<\infty.
\end{align}
因为 \(f^\varepsilon\) 是 \(C^\infty\) 的，对任意 \(x,y\in B\)，由微积分基本定理有
\begin{align*}
f^\varepsilon(x)-f^\varepsilon(y)=\int_0^1 Df^\varepsilon(y+t(x-y))\,\mathrm{d}t\cdot(x-y).
\end{align*}
两边取绝对值，并利用\eqref{eq:磨光导数一致界}式，得到
\begin{align*}
|f^\varepsilon(x)-f^\varepsilon(y)|
&\leqslant \|Df^\varepsilon\|_{L^\infty(B)}|x-y|
\leqslant \|Df\|_{L^\infty(W)}|x-y|.
\end{align*}
因此对于固定的 \(x,y\in B\)，令 \(\varepsilon\to 0^+\)，由磨光函数的性质 \(f^\varepsilon\to f\) 在 \(L^1_{\mathrm{loc}}(U)\) 中成立，故由\cref{Real Analysis-corollary:非负可测函数积分为0则几乎处处为0}知,存在子列几乎处处收敛，从而
\begin{align*}
|f(x)-f(y)|\leqslant \|Df\|_{L^\infty(W)}|x-y|.
\end{align*}
这说明 \(f|_B\) 是 Lipschitz 连续的，Lipschitz 常数不超过 \(\|Df\|_{L^\infty(W)}\)。由于闭球 \(B\subset\subset U\) 是任意的，故 \(f\) 在 \(U\) 中局部Lipschitz连续。
\end{proof}



\subsection{扩张定理和$W^{1,\infty}$的刻画}

下面我们的目标是把Sobolev空间 $W^{1,p}(U)$ 中的函数延拓为Sobolev空间 $W^{1,p}(\mathbb{R}^{n})$ 中的函数。这可能是微妙的。例如注意到，把 $u\in W^{1,p}(U)$ 在 $\mathbb{R}^{n}-U$ 中取零来延拓一般行不通，因为我们可能因此沿 $\partial U$ 造成如此恶劣的间断，以致延拓后的函数不再具有弱一阶偏导数。我们必须另想办法延拓 $u$，使之“在穿过 $\partial U$ 时保持弱导数”。

\begin{theorem}[Kirszbraun扩张定理, 实值情形]\label{theorem:实值-Kirszbraun扩张定理}
设 $A\subseteq\mathbb{R}^n$ 是任意非空子集, $f:A\to\mathbb{R}$ 是 $L$-Lipschitz 函数, 即对任意 $x,y\in A$,
\begin{align*}
|f(x)-f(y)|\leqslant L|x-y|.
\end{align*}
则存在 $L$-Lipschitz 函数 $F:\mathbb{R}^n\to\mathbb{R}$, 使得 $F|_A=f$, 且对任意 $x,y\in\mathbb{R}^n$,
\begin{align*}
|F(x)-F(y)|\leqslant L|x-y|.
\end{align*}
\end{theorem}
\begin{remark}
一般情形下, Kirszbraun 定理对向量值函数 $f:A\to\mathbb{R}^m$ 也成立, 即存在 $L$-Lipschitz 扩张 $F:\mathbb{R}^n\to\mathbb{R}^m$. 其证明需要更复杂的几何或泛函分析方法.
\end{remark}
\begin{proof}

定义
\begin{align*}
F(x)\triangleq\inf\limits_{a\in A}\left\{f(a)+L|x-a|\right\},\quad x\in\mathbb{R}^n.
\end{align*}
首先证明 $F$ 是 $L$-Lipschitz 的. 任取 $x,y\in\mathbb{R}^n$. 对任意 $a\in A$, 由三角不等式得
\begin{align*}
f(a)+L|x-a|\leqslant f(a)+L|y-a|+L|x-y|.
\end{align*}
对 $a\in A$ 取下确界, 得
\begin{align*}
F(x)\leqslant F(y)+L|x-y|.
\end{align*}
交换 $x$ 与 $y$ 的位置, 同理可得
\begin{align*}
F(y)\leqslant F(x)+L|x-y|.
\end{align*}
因此
\begin{align*}
|F(x)-F(y)|\leqslant L|x-y|.
\end{align*}

其次证明 $F|_A=f$. 任取 $b\in A$. 一方面, 对任意 $a\in A$, 由 $f$ 的 $L$-Lipschitz性得
\begin{align*}
f(a)+L|b-a|\geqslant f(b).
\end{align*}
对 $a\in A$ 取下确界, 得
\begin{align*}
F(b)=\inf\limits_{a\in A}\left\{f(a)+L|b-a|\right\}\geqslant f(b).
\end{align*}
另一方面, 在上式中取 $a=b$, 得
\begin{align*}
F(b)\leqslant f(b)+L|b-b|=f(b).
\end{align*}
故 $F(b)=f(b)$. 由 $b\in A$ 的任意性, $F|_A=f$. 证毕.
\end{proof}

\begin{theorem}[扩张定理]\label{theorem:ch05-扩张定理}
设 $1\leqslant p\leqslant\infty$，$U\subseteq \mathbb{R}^n$是有界开集且 $\partial U$ 为 $C^{1}$。选取有界开集 $V$ 使 $U\subset \subset V$。则存在有界线性算子
\begin{align}
E:W^{1,p}(U)\to W^{1,p}(\mathbb{R}^{n})
\end{align}
使对每个 $u\in W^{1,p}(U)$：
\begin{enumerate}[(i)]
\item $Eu=u$ a.e. $x\in U$；
\item $Eu$ 的支集含于 $V$,再由$V$有界知,$Eu$在$V$中有紧支集；
\item $\|Eu\|_{W^{1,p}(\mathbb{R}^{n})}\leqslant C\|u\|_{W^{1,p}(U)}$，其中常数 $C$ 只依赖于 $n,p,U,V$。
\end{enumerate}
我们称 $Eu$ 为 $u$ 到 $\mathbb{R}^{n}$ 的\textbf{扩张}。
\end{theorem}
\begin{proof}

1. 固定 \(x_0 \in \partial U\). 当 \(x_0 = 0\)，且\(\partial U\) 在 \(x_0\) 附近平坦，位于平面 \(\{x_n=0\}\) 内时，即存在 \(r>0\)，使得开球 \(B \triangleq B(0,r)\) 满足
\begin{align*}
\begin{cases}
B^{+}\triangleq B\cap\{x_{n}\geqslant 0\}\subset\overline{U},\\
B^{-}\triangleq B\cap\{x_{n}\leqslant 0\}\subset\mathbb{R}^{n}\backslash U.
\end{cases}
\end{align*}

\begin{figure}[H]
\centering
\includegraphics[width=0.55\textwidth]{ch05-half-ball.png}
\caption{边界处的半球}
\label{figure:ch05-边界半球}
\end{figure}

2. 暂时再设 $u\in C^{1}(\overline{U})$。于是定义
\begin{align}
\overline{u}(x)\triangleq
\begin{cases}
u\left( x \right) ,&		x\in B^+,\\
-3u\left( x_1,\cdots ,x_{n-1},-x_n \right) +4u\left( x_1,\cdots ,x_{n-1},-\frac{x_n}{2} \right) ,&		x\in B^-.\\
\end{cases}
\label{align:ch05-反射延拓}
\end{align}
这称为把 $u$ 从 $B^{+}$ \textbf{高阶反射}到 $B^{-}$。

3. 我们断言
\begin{align*}
\overline{u}\in C^{1}(B).
\end{align*}
首先，\(\overline{u}\) 在 \(B^+\) 上等于 \(u\)，在 \(B^-\) 上由 \eqref{align:ch05-反射延拓}式定义。由于 \(u \in C^1(\overline{U})\)，且当 \(x \in B^-\) 时，点 \((x',-x_n)\) 和 \(\left(x',-\frac{x_n}{2}\right)\) 均属于 \(B^+ \subseteq \overline{U}\)，所以 \(\overline{u}\) 在 \(B^+\) 内部和 \(B^-\) 内部均属于 \(C^1\)。因此只需验证 \(\overline{u}\) 及其一阶偏导数在边界 \(\{x_n=0\}\) 处连续。

(1) 函数的连续性：
若 \(x \in B^+\) 且 \(x_n \to 0^+\)，则由$u\in C^{1}(\overline{U})$知\(\overline{u}(x)=u(x) \to u(x',0)\)。
若 \(x \in B^-\) 且 \(x_n \to 0^-\)，则
\begin{align*}
y_n\triangleq -x_n\to 0^+,\quad z_n\triangleq -\frac{x_n}{2}\to 0^+,
\end{align*}
且$y_n,z_n\in B^+$.从而再由\eqref{align:ch05-反射延拓}式和$u\in C^{1}(\overline{U})$知
\begin{align*}
\overline{u}(x) = -3u(x',-x_n) + 4u\left(x',-\frac{x_n}{2}\right)=-3u(x',y_n) + 4u\left(x',z_n\right) \to -3u(x',0) + 4u(x',0) = u(x',0).
\end{align*}
因此 \(\overline{u}\) 在边界上连续。

(2) 一阶偏导数的连续性：
对于\(i=1,\cdots,n-1\),若\(x\in B^+\)且\(x_n \to 0^+\),则由$u\in C^{1}(\overline{U})$知
\begin{align*}
D_i \overline{u}(x) = D_i u(x) \to D_i u(x',0).
\end{align*}
若 \(x \in B^-\) 且 \(x_n \to 0^-\)，则$y_n\triangleq -x_n\to 0^+,z_n\triangleq -\frac{x_n}{2}\to 0^+$,且$y_n,z_n\in B^+$.从而再由\eqref{align:ch05-反射延拓}式和$u\in C^{1}(\overline{U})$知v
\begin{align*}
D_i \overline{u}(x) = -3D_i u(x',-x_n) + 4D_i u\left(x',-\frac{x_n}{2}\right)=-3D_i u(x',y_n) + 4D_i u\left(x',z_n\right) \to -3D_i u(x',0) + 4D_i u(x',0) = D_i u(x',0).
\end{align*}
故 \(D_i \overline{u}\) 在边界$\{x_n=0\}$上连续.

对于 \(i=n\)，若 \(x \in B^+\) 且 \(x_n \to 0^+\)，则由$u\in C^{1}(\overline{U})$知 \(D_n \overline{u}(x) = D_n u(x) \to D_n u(x',0)\)
若 \(x \in B^-\) 且 \(x_n \to 0^-\)，则
\begin{align*}
y_n\triangleq -x_n\to 0^+,\quad z_n\triangleq -\frac{x_n}{2}\to 0^+,
\end{align*}
且$y_n,z_n\in B^+$.从而再由\eqref{align:ch05-反射延拓}式和$u\in C^{1}(\overline{U})$知
\begin{align*}
D_n \overline{u}(x) = 3D_n u(x',-x_n) - 2D_n u\left(x',-\frac{x_n}{2}\right)= 3D_n u(x',y_n) - 2D_n u\left(x',z_n\right) \to 3D_n u(x',0) - 2D_n u(x',0) = D_n u(x',0).
\end{align*}
故 \(D_n \overline{u}\) 在边界$\{x_n=0\}$上连续。
综上，\(\overline{u}\) 的所有一阶偏导数在 \(B\) 上连续，因此 \(\overline{u} \in C^1(B)\)。

4. 现在证明
\begin{align*}
\|\overline{u}\|_{W^{1,p}(B)}\leqslant C\|u\|_{W^{1,p}(B^{+})},
\end{align*}
其中常数 $C$ 不依赖于 $u$。
当$1\leqslant p<\infty$时,先估计 \(L^p\) 部分. 在 \(B^+\) 上 \(\overline{u}=u\), 故
\begin{align*}
\int_{B^+}|\overline{u}|^p\,\mathrm{d}x=\int_{B^+}|u|^p\,\mathrm{d}x.
\end{align*}
在 \(B^-\) 上, 由 \eqref{align:ch05-反射延拓}式和\hyperref[Basis of Analytics-proposition:经典初等不等式]{经典初等不等式}得
\begin{align*}
|\overline{u}(x_1,\cdots,x_n)|^p
\leqslant C_p\left(3^p|u(x_1,\cdots,x_{n-1},-x_n)|^p
+4^p\left|u\left(x_1,\cdots,x_{n-1},-\frac{x_n}{2}\right)\right|^p\right),
\end{align*}
其中$C_p\triangleq 2^{p-1}$.因此
\begin{align}\label{eq::jiojiojiojioqwio2i3333-1}
\int_{B^-}|\overline{u}|^p\,\mathrm{d}x
\leqslant C_p\left(3^p\int_{B^-}|u(x_1,\cdots,x_{n-1},-x_n)|^p\,\mathrm{d}x
+4^p\int_{B^-}\left|u\left(x_1,\cdots,x_{n-1},-\frac{x_n}{2}\right)\right|^p\,\mathrm{d}x\right).
\end{align}
对第一个积分作变量替换 \(x_n\mapsto -x_n\),其他不变. 因为 \(B\) 关于 \(\{x_n=0\}\) 对称, \(B^-\) 变为 \(B^+\), 且Jacobi行列式绝对值为 \(1\), 所以
\begin{align}\label{eq::jiojiojiojioqwio2i3333-0}
\int_{B^-}|u(x_1,\cdots,x_{n-1},-x_n)|^p\,\mathrm{d}x
=\int_{B^+}|u(x_1,\cdots,x_n)|^p\,\mathrm{d}x.
\end{align}
对第二个积分作变量替换 \(x_n\mapsto -\frac{x_n}{2}\),其他不变. 此时 \(B^-\) 变为
\begin{align*}
\{x_1^2+\cdots+x_{n-1}^2+4x_n^2<r^2,\ x_n>0\},
\end{align*}
它包含于 \(B^+\), 且Jacobi行列式绝对值为 \(2\), 因此
\begin{align}\label{eq::jiojiojiojioqwio2i3333-2}
\int_{B^-}\left|u\left(x_1,\cdots,x_{n-1},-\frac{x_n}{2}\right)\right|^p\,\mathrm{d}x
=2\int_{\{x_1^2+\cdots+x_{n-1}^2+4x_n^2<r^2,\ x_n>0\}}|u(x_1,\cdots,x_n)|^p\,\mathrm{d}x
\leqslant 2\int_{B^+}|u|^p\,\mathrm{d}x.
\end{align}
合并\eqref{eq::jiojiojiojioqwio2i3333-1}\eqref{eq::jiojiojiojioqwio2i3333-0}\eqref{eq::jiojiojiojioqwio2i3333-2}式得
\begin{align*}
\int_{B^-}{\left| \overline{u} \right|^p\mathrm{d}x}\leqslant C_p\left( 3^p+2\cdot 4^p \right) \int_{B^+}{\left| u \right|^p\mathrm{d}x}.
\end{align*}
从而
\begin{align}
\left\| \overline{u} \right\| _{L^p\left( B \right)}^{p}&=\int_B{\left| \overline{u} \right|^p\mathrm{d}x}=\int_{B^+}{\left| u \right|^p\mathrm{d}x}+\int_{B^-}{\left| \overline{u} \right|^p\mathrm{d}x}
\nonumber \\
&\leqslant \int_{B^+}{\left| u \right|^p\mathrm{d}x}+C_p\left( 3^p+2\cdot 4^p \right) \int_{B^+}{\left| u \right|^p\mathrm{d}x}
\nonumber \\
&\leqslant C_p\left( 3^p+2\cdot 4^p \right) \left\| u \right\| _{L^p\left( B^+ \right)}^{p}.\label{eq:::HJiuhj3u892HFFHJIOEJO}
\end{align}
再估计一阶导数. 对 \(i=1,\cdots,n-1\), 在 \(B^-\) 上
\begin{align*}
D_i\overline{u}(x_1,\cdots,x_n)
=-3D_i u(x_1,\cdots,x_{n-1},-x_n)
+4D_i u\left(x_1,\cdots,x_{n-1},-\frac{x_n}{2}\right),
\end{align*}
利用与\eqref{eq::jiojiojiojioqwio2i3333-1}\eqref{eq::jiojiojiojioqwio2i3333-0}式同样的变量替换给出
\begin{align*}
\int_{B^-}|D_i\overline{u}|^p\,\mathrm{d}x
\leqslant C_p(3^p+2\cdot 4^p)\int_{B^+}{\left| D_iu \right|^p\mathrm{d}x}.
\end{align*}
对 \(i=n\), 在 \(B^-\) 上
\begin{align*}
D_n\overline{u}(x_1,\cdots,x_n)
=3D_n u(x_1,\cdots,x_{n-1},-x_n)
-2D_n u\left(x_1,\cdots,x_{n-1},-\frac{x_n}{2}\right),
\end{align*}
再利用与\eqref{eq::jiojiojiojioqwio2i3333-1}\eqref{eq::jiojiojiojioqwio2i3333-0}式同样的变量替换给出
\begin{align*}
\int_{B^-}|D_n\overline{u}|^p\,\mathrm{d}x
\leqslant C_p(3^p+2\cdot 2^p)\int_{B^+}{\left| D_nu \right|^p\mathrm{d}x}.
\end{align*}
在 \(B^+\) 上 \(\overline{u}=u\), 所以对每个 \(i=1,\cdots,n\),有
\begin{align}
\left\| D_i\overline{u} \right\| _{L^p\left( B \right)}^{p}&=\int_B{\left| D_i\overline{u} \right|^p\mathrm{d}x}=\int_{B^+}{\left| D_iu \right|^p\mathrm{d}x}+\int_{B^-}{\left| D_i\overline{u} \right|^p\mathrm{d}x}
\nonumber \\
&\leqslant \int_{B^+}{\left| D_iu \right|^p\mathrm{d}x}+C_p\left( 3^p+2\cdot 4^p \right) \int_{B^+}{\left| D_iu \right|^p\mathrm{d}x}
\nonumber \\
&\leqslant C_p\left( 3^p+2\cdot 4^p \right) \left\| D_iu \right\| _{L^p\left( B^+ \right)}^{p}.\label{eq:::HJiuhj3u892HFFHJIOEJO-1}
\end{align}
取$C\triangleq \left[ C_p(3^p+2\cdot 4^p) \right] ^{\frac{1}{p}}$,则由\eqref{eq:::HJiuhj3u892HFFHJIOEJO}\eqref{eq:::HJiuhj3u892HFFHJIOEJO-1}式可得
\begin{align*}
\|\overline{u}\|_{W^{1,p}(B)}^p
=\|\overline{u}\|_{L^p(B)}^p+\sum\limits_{i=1}^n\|D_i\overline{u}\|_{L^p(B)}^p\leqslant C^p\left(\|u\|_{L^p(B^+)}^p+\sum\limits_{i=1}^n\|D_i u\|_{L^p(B^+)}^p\right)=C^p\|u\|_{W^{1,p}(B^+)}^p.
\end{align*}
开 \(p\) 次方得
\begin{align*}
\|\overline{u}\|_{W^{1,p}(B)}
\leqslant C\|u\|_{W^{1,p}(B^+)}.
\end{align*}
若 \(p=\infty\), 直接取本性上界得
\begin{gather*}
\left| \overline{u} \right|\leqslant 3\left| u \right|+4\left| u \right|\leqslant 7\left\| u \right\| _{L^{\infty}\left( B^+ \right)},
\\
\left| D_i\overline{u} \right|\leqslant 3\left| D_iu \right|+4\left| D_iu \right|\leqslant 7\left\| D_iu \right\| _{L^{\infty}\left( B^+ \right)},\quad i=1,2,\cdots ,n.
\end{gather*}
取$C=7$,则有
\begin{align*}
\|\overline{u}\|_{W^{1,\infty}(B)}
\leqslant C\|u\|_{W^{1,\infty}(B^+)}.
\end{align*}

\begin{figure}[H]
\centering
\includegraphics[width=0.8\textwidth]{8e1ffee58f4e42d8a67b7ea75f26d99d.png}
\caption{拉直边界}
\label{figure:ch05-拉直边界}
\end{figure}

5.固定 \(x_0 \in \partial U\). 设 \(\nu(x_0)\) 为 \(x_0\) 处关于 \(U\) 的单位外法向量. 取$n$阶正交矩阵 \(Q \), 使得 \(Q^T\) 的最后一行为 \(\nu(x_0)^T\). 定义平移与旋转变换
\begin{align}
z &\triangleq Q(x - x_0). \label{eq:ch05-straighten-coord}
\end{align}
在此新坐标系 \(z\) 下, 点 \(x_0\) 对应 \(z = 0\), 且 \(U\) 在 \(x_0\) 处的外法向量指向 \(z_n\) 轴的正方向. 由\cref{definition:appC-Ck边界}, 存在 \(r > 0\) 和 \(C^1\) 函数 \(\gamma: \mathbb{R}^{n-1} \to \mathbb{R}\), 使得
\begin{align}
U \cap B_z(0, r) &= \{z \in B_z(0, r) \mid z_n > \gamma(z_1, \cdots, z_{n-1})\}. \label{eq:ch05-boundary-def}
\end{align}
这里 \(z' = (z_1, \cdots, z_{n-1})\). 由于 \(z_n\) 轴沿外法向, 故 \(\gamma(0) = 0\) 且 \(\nabla \gamma(0) = 0\).

定义拉直映射 \(\Phi_1: B_z(0, r) \to \mathbb{R}^n\) 为
\begin{align*}
\Phi_1(z) &\triangleq (z_1, \cdots, z_{n-1}, z_n - \gamma(z_1, \cdots, z_{n-1})).
\end{align*}
其逆映射 \(\Psi_1: \Phi_1(B_z(0, r)) \to \mathbb{R}^n\) 为
\begin{align*}
\Psi_1(y) &\triangleq (y_1, \cdots, y_{n-1}, y_n + \gamma(y_1, \cdots, y_{n-1})).
\end{align*}
直接计算可得 \(\Phi_1(\Psi_1(y)) = y\), 且 \(\Psi_1(\Phi_1(z)) = z\). 因此 \(\Phi_1\) 与 \(\Psi_1\) 互为逆映射. 由于 \(\gamma \in C^1\), 故 \(\Phi_1, \Psi_1 \in C^1\), 它们均为 \(C^1\) 微分同胚，其Jacobi矩阵分别为
\begin{align*}
D\Phi _1\left( z \right) =\left( \begin{matrix}
I_{n-1}&		0\\
-\nabla \gamma \left( z' \right) ^T&		1\\
\end{matrix} \right) ,\quad D\Psi _1\left( z \right) =\left( \begin{matrix}
I_{n-1}&		0\\
\nabla \gamma \left( z' \right) ^T&		1\\
\end{matrix} \right) .
\end{align*}
故 \(\det D\Phi_1(z) = \det D\Psi_1(y) = 1\).

现在先在 \(y\)坐标系下选取开球 \(B\). 令 \(y_0 \triangleq \Phi_1(0)=(0,\cdots,0)\). 
因为 \(\Psi_1\) 连续且 \(\Psi_1(0)=0\)，所以对任意 \(\varepsilon>0\)，存在 \(\rho_\varepsilon>0\)，使得
\begin{align*}
y\in B(0,\rho_\varepsilon)\Longrightarrow \|\Psi_1(y)\|<\varepsilon.
\end{align*}
又因为 \(Q\) 是正交矩阵，所以\(Q^T\) 保持欧氏范数，从而
\begin{align*}
W_\varepsilon\triangleq x_0+Q^T\Psi_1(B(0,\rho_\varepsilon))
\subseteq x_0+Q^TB(0,\varepsilon)
=B(x_0,\varepsilon).
\end{align*}
由于 \(x_0\in\partial U\subseteq \overline U\subseteq V\)，而 \(V\) 是开集，故存在 \(\delta>0\) 使得
\begin{align*}
B(x_0,\delta)\subseteq V.
\end{align*}
取$\varepsilon_1=\min\left\{r,\frac{\delta}{2}\right\}.$
则存在 \(\rho\triangleq \rho_{\varepsilon_1}>0\)，使得
\begin{align*}
W\triangleq W_{\varepsilon_1}\subseteq B(x_0,\varepsilon_1)\subseteq B(x_0,r)\cap B\left(x_0,\frac{\delta}{2}\right)\subseteq B(x_0,\delta)\subseteq V.
\end{align*}
故开球$B \triangleq B(y_0, \rho)$的原像
\begin{align*}
W\triangleq x_0+Q^T\Psi _1\left( B \right)
\end{align*}
包含在$B\left( x_0,r \right)$中,且$W\subset \subset V$.然后定义原坐标下的拉直映射 \(\Phi: W \to B\) 为
\begin{align*}
\Phi(x) &\triangleq \Phi_1(Q(x - x_0)),
\end{align*}
其逆映射 \(\Psi: B \to W\) 为
\begin{align*}
\Psi(y) &\triangleq x_0 + Q^T \Psi_1(y).
\end{align*}
由定义可知 \(\Phi\) 与 \(\Psi\) 互为逆映射, 且均为 \(C^1\) 微分同胚. 并且由\hyperref[theorem:向量求导基本性质]{向量求导基本性质}知
\begin{gather}
\Phi(W) = B, \quad \Psi(B) = W,\label{eq::jifjiojiojiojiomfijsfiojjsifjioejiojmwfiejfiwmfiowijijwfio-1}\\
\det D\Phi \left( x \right) =\det \left( D\Phi _1\left( x \right) Q \right) =\det D\Phi _1\left( x \right) \cdot \det Q=1,\label{eq::jifjiojiojiojiomfijsfiojjsifjioejiojmwfiejfiwmfiowijijwfio-2}
\\
\det D\Psi \left( x \right) =\det \left( Q^TD\Psi _1\left( x \right) \right) =\det Q^T\cdot \det D\Phi _1\left( x \right) =1.\label{eq::jifjiojiojiojiomfijsfiojjsifjioejiojmwfiejfiwmfiowijijwfio-3}
\end{gather}

下面验证 \(\Phi\) 将边界拉直. 若 \(x \in W \cap U\), 则 \(z = Q(x - x_0) \in B_z(0, r)\), 且 \(z_n > \gamma(z')\). 于是
\begin{align*}
\Phi_n(x) &= (\Phi_1(z))_n = z_n - \gamma(z') > 0,
\end{align*}
即 \(\Phi(x) \in B^+\). 因此 \(\Phi(W \cap U) \subseteq B^+\). 反之, 若 \(y \in B^+\), 则 \(y_n > 0\). 令 \(x = \Psi(y)\), 则 \(z = Q(x - x_0) = \Psi_1(y)\). 于是
\begin{align*}
z_n &= y_n + \gamma(y') > \gamma(y') = \gamma(z'),
\end{align*}
故 \(x \in W \cap U\). 因此 \(\Phi(W \cap U) = B^+\). 若 \(x \in W \cap \partial U\), 则 \(z_n = \gamma(z')\), 于是 \(\Phi_n(x) = 0\), 即 \(\Phi(x) \in B \cap \{y_n = 0\}\). 因此 \(\Phi(W \cap \partial U) = B \cap \{y_n = 0\}\).

下面在 \(B\) 中作反射延拓. 先设 \(u \in C^1(\overline{U})\). 定义
\begin{align*}
u'(y) \triangleq u(\Psi(y)), \quad y \in B^+.
\end{align*}
对 \(u'\) 应用步骤 1--4 中的高阶反射延拓, 得到 \(\overline{u}' \in C^1(B)\), 满足
\begin{gather}
\overline{u}'(y) = u'(y), \quad y \in B^+, \label{eq:ch05-ext-in-ball-1} \\
\|\overline{u}'\|_{W^{1,p}(B)} \leqslant C \|u'\|_{W^{1,p}(B^+)}. \label{eq:ch05-ext-in-ball-2}
\end{gather}
由于 \(u'\) 在 \(B^+\) 中具有紧支集, 故 \(\overline{u}'\) 在 \(B\) 中具有紧支集.

定义 \(\overline{u}(x) \triangleq \overline{u}'(\Phi(x))\), \(x \in W\). 则 \(\overline{u} \in C^1(W)\), 且 \(\overline{u}\) 在 \(W\) 中具有紧支集. 若 \(x \in W \cap U\), 则 \(\Phi(x) \in B^+\), 于是
\begin{align*}
\overline{u}(x) &= \overline{u}'(\Phi(x)) = u'(\Phi(x)) = u(\Psi(\Phi(x))) = u(x).
\end{align*}
故 \(\overline{u}|_{W \cap U} = u\). 由变量替换 \(y = \Phi(x)\), \(x = \Psi(y)\) 及\eqref{eq::jifjiojiojiojiomfijsfiojjsifjioejiojmwfiejfiwmfiowijijwfio-1}\eqref{eq::jifjiojiojiojiomfijsfiojjsifjioejiojmwfiejfiwmfiowijijwfio-2}\eqref{eq::jifjiojiojiojiomfijsfiojjsifjioejiojmwfiejfiwmfiowijijwfio-3}式, 得
\begin{gather}
\int_W |\overline{u}(x)|^p \mathrm{d}x = \int_B |\overline{u}'(y)|^p |\det D\Psi(y)| \mathrm{d}y = \int_B |\overline{u}'(y)|^p \mathrm{d}y,\label{eq::jijioqwio93283r3r23}\\
\int_{B^+} |u'(y)|^p \mathrm{d}y = \int_{W \cap U} |u(x)|^p |\det D\Phi(x)| \mathrm{d}x = \int_{W \cap U} |u(x)|^p \mathrm{d}x.\label{eq::jijioqwio93283r3r23-111}
\end{gather}
由链式法则得
\begin{align}
D\overline{u}(x) = D\overline{u}'(\Phi(x)) D\Phi(x). \label{eq:ch05-chain-rule-uv}
\end{align}
对$\forall A\in \mathbb{R}^{n\times n}$,定义矩阵范数$\left\| A \right\| \triangleq \underset{v\in \mathbb{R} ^n\backslash \left\{ 0 \right\}}{\mathrm{sup}}\frac{\left\| Av \right\|}{\left\| v \right\|}$.下面证明 \(\|D\Phi(x)\|\) 在 \(\overline{W}\) 上关于 \(x\) 有界。因为 \(\Phi \in C^1\)，所以 \(D\Phi(x)\) 的每个元素 \(\frac{\partial \Phi_i}{\partial x_j}(x)\) 在 \(\overline{W}\) 上连续，从而矩阵值函数 \(x \mapsto D\Phi(x)\) 连续。
又由Cauchy-Schwarz不等式，对任意单位向量 \(v \in \mathbb{R}^n\)，有
\begin{align*}
\|Av\|^2 &= \sum_{i=1}^n \left( \sum_{j=1}^n a_{ij} v_j \right)^2 \leqslant \sum_{i=1}^n \left( \sum_{j=1}^n a_{ij}^2 \right) \left( \sum_{j=1}^n v_j^2 \right) = \sum_{i,j=1}^n a_{ij}^2.
\end{align*}
因此
\begin{align*}
\|A\| &\leqslant \left( \sum_{i,j=1}^n a_{ij}^2 \right)^{\frac{1}{2}}.
\end{align*}
于是当 \(x_k \to x\) 时，结合\cref{Functional Analysis-theorem:有界线性算子全体构成的空间完备}有
\begin{align*}
\bigl| \|D\Phi(x_k)\| - \|D\Phi(x)\| \bigr| \leqslant \|D\Phi(x_k) - D\Phi(x)\| &\leqslant \left( \sum_{i,j=1}^n \left( \frac{\partial \Phi_i}{\partial x_j}(x_k) - \frac{\partial \Phi_i}{\partial x_j}(x) \right)^2 \right)^{\frac{1}{2}} \to 0.
\end{align*}
故 \(x \mapsto \|D\Phi(x)\|\) 在 \(\overline{W}\) 上连续。由于 \(\overline{W}\) 是紧集，连续函数在紧集上必有最大值，记
\begin{align*}
C_1 &\triangleq \max_{x \in \overline{W}} \|D\Phi(x)\|,
\end{align*}
于是对一切 \(x \in \overline{W}\)，有 \(\|D\Phi(x)\| \leqslant C_1\)。由\cref{Functional Analysis-proposition:有界线性算子的基本不等式}，对任意 \(v \in \mathbb{R}^n\)，有
\begin{align*}
|v D\Phi(x)| \leqslant \|D\Phi(x)\| |v|.
\end{align*}
取 \(v = D\overline{u}'(\Phi(x))\)，并结合 \eqref{eq:ch05-chain-rule-uv}，得到
\begin{align}\label{eq::jioJFIOWJIOIJOjiojfwio-1}
|D\overline{u}(x)| = |D\overline{u}'(\Phi(x)) D\Phi(x)| \leqslant \|D\Phi(x)\| |D\overline{u}'(\Phi(x))| \leqslant C_1 |D\overline{u}'(\Phi(x))|.
\end{align}
同理可得,存在$C_2>0$,使得
\begin{align}\label{eq::jioJFIOWJIOIJOjiojfwio-2}
|Du'(x)| = |Du'(\Psi(x)) D\Psi(x)| \leqslant \|D\Psi(x)\| |Du'(\Psi(x))| \leqslant C_2 |Du'(\Psi(x))|.
\end{align}
利用\eqref{eq::jioJFIOWJIOIJOjiojfwio-1}\eqref{eq::jioJFIOWJIOIJOjiojfwio-2}式,再由变量替换 \(y = \Phi(x)\), \(x = \Psi(y)\) 及\eqref{eq::jifjiojiojiojiomfijsfiojjsifjioejiojmwfiejfiwmfiowijijwfio-1}\eqref{eq::jifjiojiojiojiomfijsfiojjsifjioejiojmwfiejfiwmfiowijijwfio-2}\eqref{eq::jifjiojiojiojiomfijsfiojjsifjioejiojmwfiejfiwmfiowijijwfio-3}式, 得
\begin{gather}
\int_W |D\overline{u}(x)|^p \mathrm{d}x \leqslant C_1^p \int_W |D\overline{u}'(\Phi(x))|^p \mathrm{d}x = C_1^p \int_B |D\overline{u}'(y)|^p |\det D\Phi(y)| \mathrm{d}y=C_1^p \int_B |D\overline{u}'(y)|^p \mathrm{d}y,\label{eq::jijioqwio93283r3r23-1}
\\
\int_{B^+}{\left| Du' \left( x \right) \right|^p\mathrm{d}x}\leqslant C_{2}^{p}\int_{B^+}{\left| Du' \left( \Psi \left( x \right) \right) \right|^p\mathrm{d}x}=C_{2}^{p}\int_{W\cap U}{\left| Du' \left( y \right) \right|^p |\det D\Psi(y)|\mathrm{d}y}=C_{2}^{p}\int_{W\cap U}{\left| Du' \left( y \right) \right|^p \mathrm{d}y}.\label{eq::jijioqwio93283r3r23-2}
\end{gather}
结合\eqref{eq::jijioqwio93283r3r23}\eqref{eq::jijioqwio93283r3r23-111}\eqref{eq::jijioqwio93283r3r23-1}\eqref{eq::jijioqwio93283r3r23-2}式可得
\begin{align}
\left\| \overline{u} \right\| _{W^{1,p}\left( W \right)}&=\left( \int_W{\left| \overline{u}\left( x \right) \right|^p\mathrm{d}x}+\int_W{\left| D\overline{u}\left( x \right) \right|^p\mathrm{d}x} \right) ^{\frac{1}{p}}\nonumber
\\
&\leqslant \left( \int_B{\left| \overline{u}\left( y \right) \right|^p\mathrm{d}y}+C_{1}^{p}\int_B{\left| D\overline{u}\left( y \right) \right|^p\mathrm{d}y} \right) ^{\frac{1}{p}}\nonumber\\
&\leqslant C_3\left\| \overline{u'} \right\| _{W^{1,p}\left( B \right)},
\label{eq:ch05-local-est-1}
\end{align}
\begin{align}
\left\| u^{\prime} \right\| _{W^{1,p}\left( B^+ \right)}&=\left( \int_{B^+}{\left| u^{\prime}\left( x \right) \right|^p\mathrm{d}x}+\int_{B^+}{\left| Du^{\prime}\left( x \right) \right|^p\mathrm{d}x} \right) ^{\frac{1}{p}}
\nonumber\\
&\leqslant \left( \int_{W\cap U}{\left| u^{\prime}\left( y \right) \right|^p\mathrm{d}y}+C_{2}^{p}\int_{W\cap U}{\left| Du^{\prime}\left( y \right) \right|^p\mathrm{d}y} \right) ^{\frac{1}{p}}
\nonumber\\
&\leqslant C_3\left\| \overline{u'} \right\| _{W^{1,p}\left( W\cap U \right)}\leqslant C_3\left\| \overline{u'} \right\| _{W^{1,p}\left( U \right)},\label{eq:ch05-local-est-2}
\end{align}
其中$C_3=\max\{C_1,C_2,1\}$.再记$C'\triangleq CC_3^2$,将 \eqref{eq:ch05-ext-in-ball-2}\eqref{eq:ch05-local-est-1} \eqref{eq:ch05-local-est-2}式合并, 得到
\begin{align}
\left\| \overline{u} \right\| _{W^{1,p}\left( W \right)}\leqslant C_3\left\| \overline{u^{\prime}} \right\| _{W^{1,p}\left( B \right)}\leqslant CC_3\left\| u^{\prime} \right\| _{W^{1,p}\left( B^+ \right)}\leqslant CC_{3}^{2}\left\| \overline{u^{\prime}} \right\| _{W^{1,p}\left( U \right)}=C^{\prime}\left\| \overline{u^{\prime}} \right\| _{W^{1,p}\left( U \right)}.\label{eq:ch05-local-est-final}
\end{align}

若 \(p=\infty\)，先说明 \(L^\infty\) 范数在 \(C^1\) 微分同胚下不变。由\cref{Real Analysis-corollary:Lebesgue积分基本定理}和\subcref{Real Analysis-proposition:绝对连续函数零测集的像仍零测}{Real Analysis-proposition:绝对连续函数零测集的像仍零测-1}可知，\(\Phi\) 与 \(\Psi\) 均把零测集映为零测集.设 \(f:B\to\mathbb{R}\),$\,g:W\cap U\to \mathbb{R}$均可测,下面证明
\begin{align}\label{eq::jiowiowijoijoijijojsiosioios3829237489}
\|f\circ\Phi\|_{L^\infty(W)}=\|f\|_{L^\infty(B)},\qquad
\|g\circ\Psi\|_{L^\infty(B^+)}=\|g\|_{L^\infty(W\cap U)}.
\end{align}
若 \(\|f\|_{L^\infty(B)}=\infty\)，则对任意 \(M>0\)，集合
\begin{align*}
N_M\triangleq\{y\in B:|f(y)|>M\}
\end{align*}
不是零测集。由 \(\Phi\) 把零测集映为零测集，故若 \(\Psi(N_M)\) 为零测集，则
\begin{align*}
N_M=\Phi(\Psi(N_M))
\end{align*}
为零测集，矛盾！因此 \(\Psi(N_M)\) 不是零测集。而
\begin{align*}
\{x\in W:|f(\Phi(x))|>M\}=\Phi^{-1}(N_M)=\Psi(N_M),
\end{align*}
故对任意 \(M>0\)，该集合不是零测集，所以
\begin{align*}
\|f\circ\Phi\|_{L^\infty(W)}=\infty.
\end{align*}
等号成立。

若 \(\|f\|_{L^\infty(B)}<\infty\)，记$M_0\triangleq\|f\|_{L^\infty(B)},$
则 \(|f|\leqslant M_0\) 几乎处处于 \(B\)。令
\begin{align*}
N\triangleq\{y\in B:|f(y)|>M_0\},
\end{align*}
则 \(N\) 为零测集。由 \(\Psi\) 把零测集映为零测集知\(\Psi(N)\) 为零测集。而
\begin{align*}
\{x\in W:|f(\Phi(x))|>M_0\}=\Phi^{-1}(N)=\Psi(N),
\end{align*}
故 \(|f\circ\Phi|\leqslant M_0\) 几乎处处于 \(W\)。因此
\begin{align}\label{eq::jiojfwiojiojwio3i2900384}
\|f\circ\Phi\|_{L^\infty(W)}\leqslant M_0=\|f\|_{L^\infty(B)}.
\end{align}
反过来，记$M_1\triangleq\|f\circ\Phi\|_{L^\infty(W)},$
则 \(|f\circ\Phi|\leqslant M_1\) 几乎处处于 \(W\)。令
\begin{align*}
N'\triangleq\{x\in W:|f(\Phi(x))|>M_1\}.
\end{align*}
则 \(N'\) 为零测集。由 \(\Phi\) 把零测集映为零测集知\(\Phi(N')\) 为零测集。而
\begin{align*}
\Phi(N')=\{y\in B:|f(y)|>M_1\},
\end{align*}
故 \(|f|\leqslant M_1\) 几乎处处于 \(B\)。因此
\begin{align}\label{eq::jiojfwiojiojwio3i2900382}
\|f\|_{L^\infty(B)}\leqslant M_1=\|f\circ\Phi\|_{L^\infty(W)}.
\end{align}
合并\eqref{eq::jiojfwiojiojwio3i2900384}\eqref{eq::jiojfwiojiojwio3i2900382}式，得
\begin{align*}
\|f\circ\Phi\|_{L^\infty(W)}=\|f\|_{L^\infty(B)}.
\end{align*}
把 \(f\) 换成 \(g\)，把 \(\Phi\) 换成 \(\Psi\)，并将 \(W\)换成$B^+$,\(B\)换成$W\cap U$,同理可得
\begin{align*}
\|g\circ\Psi\|_{L^\infty(B^+)}=\|g\|_{L^\infty(W\cap U)}.
\end{align*}

现在一方面,由 \(\overline u(x)=\overline u'(\Phi(x)),x\in W\)和\eqref{eq::jiowiowijoijoijijojsiosioios3829237489}式，得到
\begin{align*}
\|\overline u\|_{L^\infty(W)}
=\|\overline u'\circ\Phi\|_{L^\infty(W)}
=\|\overline u'\|_{L^\infty(B)}.
\end{align*}
再结合\eqref{eq::jioJFIOWJIOIJOjiojfwio-1}\eqref{eq::jiowiowijoijoijijojsiosioios3829237489}式可得
\begin{align*}
\|D\overline u\|_{L^\infty(W)}
\leqslant C_1\|D\overline u'\circ\Phi\|_{L^\infty(W)}
=C_1\|D\overline u'\|_{L^\infty(B)}.
\end{align*}
于是
\begin{align}\label{eq::jiwjiofijoswjiowiow2390jw89023890}
\left\| \overline{u} \right\| _{W^{1,\infty}\left( W \right)}=\left\| \overline{u} \right\| _{L^{\infty}\left( W \right)}+\left\| D\overline{u} \right\| _{L^{\infty}\left( W \right)}\leqslant \left\| \overline{u}^{\prime} \right\| _{L^{\infty}\left( B \right)}+C_1\left\| D\overline{u}^{\prime} \right\| _{L^{\infty}\left( B \right)}\leqslant C_3\left\| \overline{u}^{\prime} \right\| _{W^{1,\infty}\left( B \right)}.
\end{align}

另一方面，由 \(u'(y)=u(\Psi(y))\)，\(y\in B^+\)和\eqref{eq::jiowiowijoijoijijojsiosioios3829237489}式有
\begin{align}
\left\| u' \right\| _{L^{\infty}\left( B^+ \right)}=\left\| u\circ \Psi \right\| _{L^{\infty}\left( B^+ \right)}=\left\| u \right\| _{L^{\infty}\left( W\cap U \right)}
\end{align}
再结合\eqref{eq::jioJFIOWJIOIJOjiojfwio-2}\eqref{eq::jiowiowijoijoijijojsiosioios3829237489}式可得
\begin{align*}
\|Du'\|_{L^\infty(B^+)}
\leqslant C_2\|Du\circ\Psi\|_{L^\infty(B^+)}
=C_2\|Du\|_{L^\infty(W\cap U)}.
\end{align*}
从而
\begin{align}
\left\| u^{\prime} \right\| _{W^{1,\infty}\left( B^+ \right)}&=\left\| u^{\prime} \right\| _{L^{\infty}\left( B^+ \right)}+\left\| Du^{\prime} \right\| _{L^{\infty}\left( B^+ \right)}\leqslant \left\| u \right\| _{L^{\infty}\left( W\cap U \right)}+C_2\left\| Du \right\| _{L^{\infty}\left( W\cap U \right)}
\leqslant C_3\left\| u \right\| _{W^{1,\infty}\left( W\cap U \right)}.\label{eq::jiwjiofijoswjiowiow2390jw89023890-1}
\end{align}
将\eqref{eq:ch05-ext-in-ball-2}\eqref{eq::jiwjiofijoswjiowiow2390jw89023890}\eqref{eq::jiwjiofijoswjiowiow2390jw89023890-1}式合并，得到
\begin{align*}
\|\overline u\|_{W^{1,\infty}(W)}
\leqslant C_3\|\overline u'\|_{W^{1,\infty}(B)}
\leqslant C C_3\|u'\|_{W^{1,\infty}(B^+)}
\leqslant C C_3^2\|u\|_{W^{1,\infty}(W\cap U)}\leqslant C C_3^2\|u\|_{W^{1,\infty}(U)}=C'\|u\|_{W^{1,\infty}(U)}.
\end{align*}

综上，对任意给定的 \(u\in C^1(\overline{U})\) 以及任意 \(x_0\in\partial U\)，存在开集\(W\)，满足 \(x_0\in W\subset\subset V\)，并存在函数 \(\overline{u}\in C^1(W)\) 和只依赖于 \(p,U,V\) 的常数 \(C'>0\)，使得
\begin{gather}
\overline{u}=u\quad \text{a.e.}\, x\in W\cap U,\nonumber \\
\|\overline{u}\|_{W^{1,p}(W)}
\leqslant C'\|u\|_{W^{1,p}(U)},\quad 1\leqslant p\leqslant \infty.\label{eq::eEREFFwww}
\end{gather}

6. 由$\partial U$ 紧和第5步知，存在有限个点 $x_{i}^{0}\in\partial U$和包含$x_i^0$的开集 $W_{i}\subset \subset V$ 以及$u$ 到 $W_{i}$ 的延拓 $\overline{u}_{i}$（$i=1,\cdots,N$），使 $\partial U\subset\bigcup_{i=1}^{N}W_{i}$,且
\begin{align}
\|\overline{u}_i\|_{W^{1,p}(W_i)}
\leqslant C'\|u_i\|_{W^{1,p}(U)},\quad i=1,2,\cdots,N,\,1\leqslant p\leqslant \infty.\label{eq::eEREFFwww1}
\end{align}
因为 \(\partial U\) 紧且 \(\partial U\subseteq\bigcup_{i=1}^N W_i\)，其中每个 \(W_i\) 是开集。考虑
\begin{align}\label{eq:ioujiwjioKKKI}
K\triangleq U\setminus\bigcup_{i=1}^N W_i.
\end{align}
先证明 \(\overline{K}\subseteq U\)。若存在 \(x\in\overline{K}\cap\partial U\)，则由 \(\partial U\subseteq\bigcup_{i=1}^N W_i\) 知存在 \(j\in\{1,\cdots,N\}\) 使得 \(x\in W_j\)。因为 \(W_j\) 是开集，所以存在 \(\delta>0\) 使得 \(B(x,\delta)\subseteq W_j\)。又 \(x\in\overline{K}\)，故 \(B(x,\delta)\cap K\neq\varnothing\)。这与$K\subseteq U\setminus\bigcup_{i=1}^N W_i\subseteq \mathbb{R}^n\setminus W_j$矛盾!
因此 \(\overline{K}\cap\partial U=\varnothing\)。又由\eqref{eq:ioujiwjioKKKI}式和\subcref{Real Analysis-theorem:边界和开集基本性质}{Real Analysis-theorem:边界和开集基本性质-5}知\(\overline{K}\subseteq \overline{U}=U\cup \partial U\),故$\overline{K}\subseteq U$.又 \(\overline{K}\) 是闭集且 \(\overline{K}\subseteq\overline{U}\)，而 \(\overline{U}\) 紧，故 \(\overline{K}\) 紧。

因为 \(\overline{K}\) 紧且 \(\overline{K}\subseteq U\)，而 \(U\) 是有界开集，所以由\subcref{Real Analysis-proposition:点集间的距离的性质}{Real Analysis-proposition:点集间的距离的性质-3}知$\mathrm{dist}(\overline{K},\mathbb{R}^n\setminus U)\geqslant \mathrm{dist}(\overline{K},\partial U)>0$.令$\eta\triangleq\frac{1}{2}\mathrm{dist}(\overline{K},\mathbb{R}^n\setminus U)>0$,并定义
\begin{align*}
W_0\triangleq\{x\in\mathbb{R}^n:\mathrm{dist}(x,\overline{K})<\eta\}.
\end{align*}
则由距离函数连续知\(W_0\) 是开集，且
\begin{align*}
K\subseteq \overline{K}\subseteq \overline{W_0}\subseteq\{x\in\mathbb{R}^n:\mathrm{dist}(x,\overline{K})\leqslant\eta\}\subseteq U.
\end{align*}
因此 \(W_0\subset\subset U\)。

最后验证 \(U\subseteq W_0\cup\bigcup_{i=1}^N W_i\)。任取 \(x\in U\)。若 \(x\in\bigcup_{i=1}^N W_i\)，则已属于右边。若 \(x\notin\bigcup_{i=1}^N W_i\)，则 \(x\in K\subseteq\overline{K}\subseteq W_0\)。故无论如何都有 \(x\in W_0\cup\bigcup_{i=1}^N W_i\)。因此
\begin{align*}
U\subseteq W_0\cup\bigcup_{i=1}^N W_i=\bigcup_{i=0}^N W_i.
\end{align*}
由\hyperref[lemma:ch05-光滑单位分解-p7R3]{光滑单位分解定理}知,
对$i=0,1,\cdots,N$,存在 $\zeta_i\in C_c^\infty(W_i)$, 使得
\begin{align*}
0\leqslant\zeta_i\leqslant1,\quad
\sum\limits_{i=0}^{N}\zeta_i(x)=1,\quad\forall x\in U.
\end{align*}
记$\overline{u}\triangleq\sum_{i=0}^{N}\zeta_{i}\overline{u}_{i}$，其中 $\overline{u}_{0}=u$。
因为 \(\zeta_i\in C_c^\infty(W_i)\)，所以 \(\zeta_i\) 具有紧支集。设
\[
K_i\triangleq \operatorname{spt}\zeta_i\subset\subset W_i.
\]
由于 \(\zeta_i\) 是连续函数，而 \(K_i\) 是紧集，所以 \(\zeta_i\) 在 \(K_i\) 上有界，即
\[
\|\zeta_i\|_{L^\infty(W_i)}=\sup_{x\in W_i}|\zeta_i(x)|=\max_{x\in K_i}|\zeta_i(x)|<\infty.
\]
同理，对每个 \(j=1,\cdots,n\)，偏导数 \(D_j\zeta_i\) 也是连续的。由于 \(\zeta_i\) 在 \(W_i\setminus K_i\) 上恒为零，所以 \(D_j\zeta_i\) 也在 \(W_i\setminus K_i\) 上恒为零，因此
\[
\operatorname{spt}D_j\zeta_i\subseteq K_i.
\]
于是 \(D_j\zeta_i\) 也在紧集 \(K_i\) 上连续，从而有界，即
\[
\|D_j\zeta_i\|_{L^\infty(W_i)}=\max_{x\in K_i}|D_j\zeta_i(x)|<\infty.
\]
因此
\[
\|D\zeta_i\|_{L^\infty(W_i)}\triangleq \max_{1\leqslant j\leqslant n}\|D_j\zeta_i\|_{L^\infty(W_i)}<\infty.
\]
所以
\[
M_i\triangleq\|\zeta_i\|_{L^\infty(W_i)}^p+\|D\zeta_i\|_{L^\infty(W_i)}^p<\infty,\quad M_i'\triangleq \|\zeta_i\|_{L^\infty(W_i)}+\|D\zeta_i\|_{L^\infty(W_i)}<\infty.
\]
对每个\(i=0,1,\cdots,N\)，由\hyperref[theorem:ch05-弱导数性质-4]{Leibniz公式}得
\begin{align*}
D(\zeta_i \overline{u}_i) = \overline{u}_i D\zeta_i + \zeta_i D\overline{u}_i.
\end{align*}
因此当 \(1\leqslant p<\infty\) 时，有
\begin{align*}
\left\| \zeta _i\overline{u}_i \right\| _{W^{1,p}\left( W_i \right)}&=\left( \left\| \zeta _i\overline{u}_i \right\| _{L^p\left( W_i \right)}^{p}+\left\| \overline{u}_iD\zeta _i \right\| _{L^p\left( W_i \right)}^{p}+\left\| \zeta _iD\overline{u}_i \right\| _{L^p\left( W_i \right)}^{p} \right) ^{\frac{1}{p}}
\\
&\leqslant \left( \left\| \zeta _i \right\| _{L^{\infty}\left( W_i \right)}^{p}\left\| \overline{u}_i \right\| _{L^p\left( W_i \right)}^{p}+\left\| \zeta _i \right\| _{L^{\infty}\left( W_i \right)}^{p}\left\| D\overline{u}_i \right\| _{L^p\left( W_i \right)}^{p}+\left\| D\zeta _i \right\| _{L^{\infty}\left( W_i \right)}^{p}\left\| \overline{u}_i \right\| _{L^p\left( W_i \right)}^{p} \right) ^{\frac{1}{p}}
\\
&\leqslant \left[ \left( \left\| \zeta _i \right\| _{L^{\infty}\left( W_i \right)}^p+\left\| D\zeta _i \right\| _{L^{\infty}\left( W_i \right)}^{p} \right) \left\| \overline{u}_i \right\| _{L^p\left( W_i \right)}^{p}+\left\| \zeta _i \right\| _{L^{\infty}\left( W_i \right)}^{p}\left\| D\overline{u}_i \right\| _{L^p\left( W_i \right)}^{p} \right] ^{\frac{1}{p}}
\\
&\leqslant \left( \left\| \zeta _i \right\| _{L^{\infty}\left( W_i \right)}^p+\left\| D\zeta _i \right\| _{L^{\infty}\left( W_i \right)}^{p} \right) ^{\frac{1}{p}}\left( \left\| \overline{u}_i \right\| _{L^p\left( W_i \right)}^{p}+\left\| D\overline{u}_i \right\| _{L^p\left( W_i \right)}^{p} \right) ^{\frac{1}{p}}
\\
&\leqslant M_{i}^{\frac{1}{p}}\left\| \overline{u}_i \right\| _{W^{1,p}\left( W_i \right)}.
\end{align*}
从而
\begin{align*}
\|\overline{u}\|_{W^{1,p}(\mathbb{R}^n)}
&\leqslant \sum\limits_{i=0}^N \|\zeta_i \overline{u}_i\|_{W^{1,p}(W_i)}
\leqslant \sum\limits_{i=0}^N M_i^{\frac{1}{p}} \|\overline{u}_i\|_{W^{1,p}(W_i)}\\
&\leqslant M_0^{\frac{1}{p}} \|u\|_{W^{1,p}(U)} + \sum\limits_{i=1}^N M_i^{\frac{1}{p}} C' \|u\|_{W^{1,p}(U)}.
\end{align*}
于是取$C'' \triangleq M_0^{\frac{1}{p}} + C' \sum\limits_{i=1}^N M_i^{\frac{1}{p}},$则
\begin{align*}
\|\overline{u}\|_{W^{1,p}(\mathbb{R}^n)}\leqslant C''\|u\|_{W^{1,p}(U)},
\end{align*}
其中常数 \(C''\) 只依赖于 \(U,V,p,n\) 以及单位分解的选取,而单位分解的选取只依赖于\(U,V,p,n\),故常数 \(C''\) 只依赖于 \(U,V,p,n\).

当$p=\infty$时,有
\begin{align*}
\left\| \zeta _i\overline{u}_i \right\| _{W^{1,\infty}\left( W_i \right)}&=\underset{x\in W_i}{\mathrm{ess}\,\mathrm{sup}}\left( \zeta _i\overline{u}_i \right) +\underset{x\in W_i}{\mathrm{ess}\,\mathrm{sup}}\left( \overline{u}_iD\zeta _i \right) +\underset{x\in W_i}{\mathrm{ess}\,\mathrm{sup}}\left( \zeta _iD\overline{u}_i \right) 
\\
&\leqslant \left\| \zeta _i \right\| _{L^{\infty}\left( W_i \right)}\underset{x\in W_i}{\mathrm{ess}\,\mathrm{sup}}\left( \overline{u}_i \right) +\left\| \zeta _i \right\| _{L^{\infty}\left( W_i \right)}\underset{x\in W_i}{\mathrm{ess}\,\mathrm{sup}}\left( D\overline{u}_i \right) +\left\| D\zeta _i \right\| _{L^{\infty}\left( W_i \right)}\underset{x\in W_i}{\mathrm{ess}\,\mathrm{sup}}\left( \overline{u}_i \right) 
\\
&\leqslant \left( \left\| \zeta _i \right\| _{L^{\infty}\left( W_i \right)}+\left\| D\zeta _i \right\| _{L^{\infty}\left( W_i \right)} \right) \underset{x\in W_i}{\mathrm{ess}\,\mathrm{sup}}\left( \overline{u}_i \right) +\left\| \zeta _i \right\| _{L^{\infty}\left( W_i \right)}\underset{x\in W_i}{\mathrm{ess}\,\mathrm{sup}}\left( D\overline{u}_i \right) 
\\
&\leqslant \left( \left\| \zeta _i \right\| _{L^{\infty}\left( W_i \right)}+\left\| D\zeta _i \right\| _{L^{\infty}\left( W_i \right)} \right) \left( \underset{x\in W_i}{\mathrm{ess}\,\mathrm{sup}}\left( \overline{u}_i \right) +\underset{x\in W_i}{\mathrm{ess}\,\mathrm{sup}}\left( D\overline{u}_i \right) \right) 
\\
&\leqslant M_{i}^{\prime}\left\| \overline{u}_i \right\| _{W^{1,\infty}\left( W_i \right)}.
\end{align*}
又因为$W_i\subset \subset V$，所以
\begin{align*}
\mathrm{spt}\,\overline{u}\subseteq \bigcup_{i=0}^N{\mathrm{spt}\,\zeta _i}\subseteq \bigcup_{i=0}^N{W_i}\subset \subset V.
\end{align*}
从而
\begin{align*}
\left\| \overline{u} \right\| _{W^{1,\infty}\left( \mathbb{R} ^n \right)}\leqslant \sum_{i=0}^N{\left\| \zeta _i\overline{u}_i \right\| _{W^{1,\infty}\left( W_i \right)}}\leqslant \sum_{i=0}^N{M_{i}^{\prime}}\left\| \overline{u}_i \right\| _{W^{1,\infty}\left( W_i \right)}
\\
\leqslant M_{0}^{\prime}\left\| u \right\| _{W^{1,\infty}\left( U \right)}+\sum_{i=1}^N{M_{i}^{\prime}C\prime}\left\| u \right\| _{W^{1,\infty}\left( U \right)}.
\end{align*}
于是取$C'' \triangleq M_0' + C' \sum\limits_{i=1}^N M_i',$则
\begin{align*}
\|\overline{u}\|_{W^{1,p}(\mathbb{R}^n)}\leqslant C''\|u\|_{W^{1,p}(U)},
\end{align*}
其中常数 \(C''\) 只依赖于 \(U,V,p,n\) 以及单位分解的选取,而单位分解的选取只依赖于\(U,V,p,n\),故常数 \(C''\) 只依赖于 \(U,V,p,n\).综上,对$1\leqslant p\leqslant \infty$,都存在只依赖于 \(U,V,p,n\)的常数$C''>0$,使得
\begin{align}\label{eq::fjiwjio233--12}
\|\overline{u}\|_{W^{1,p}(\mathbb{R}^n)}\leqslant C''\|u\|_{W^{1,p}(U)},
\end{align}

7. 对于 \(u\in C^{\infty}(\overline{U})\)，由1--6步构造的 \(\overline{u}\) 定义$E_0 u\triangleq \overline{u},$
则由1-6步的构造都是线性的可知\(E_0:C^{\infty}(\overline{U})\to W^{1,p}(\mathbb{R}^n)\) 是线性的，且由\eqref{eq::fjiwjio233--12}式有
\begin{align}\label{eqJIOFJIOWJIO323u89rfu3}
\|E_0 u\|_{W^{1,p}(\mathbb{R}^n)}\leqslant C''\|u\|_{W^{1,p}(U)}.
\end{align}

现设 \(1\leqslant p<\infty\)，\(u\in W^{1,p}(U)\)。由\cref{theorem:ch05-直到边界光滑的函数整体逼近定理}，可选 \(u_m\in C^{\infty}(\overline{U})\) 在 \(W^{1,p}(U)\) 中收敛于 \(u\)。由\eqref{eqJIOFJIOWJIO323u89rfu3}式与 \(E_0\) 的线性性可得
\begin{align*}
\|E_0 u_m-E_0 u_l\|_{W^{1,p}(\mathbb{R}^n)}
&\leqslant C''\|u_m-u_l\|_{W^{1,p}(U)}.
\end{align*}
故 \(\{E_0 u_m\}_{m=1}^{\infty}\) 是 \(W^{1,p}(\mathbb{R}^n)\) 中的 Cauchy 列。由\cref{theorem:ch05-WkpBanach}，\(W^{1,p}(\mathbb{R}^n)\) 是 Banach 空间，故存在 \(\overline{u}\in W^{1,p}(\mathbb{R}^n)\) 使得
\begin{align*}
E_0 u_m\to \overline{u}\quad \text{于 } W^{1,p}(\mathbb{R}^n).
\end{align*}
定义$E u\triangleq \overline{u}.$
下证该$E$是良定义的。假设 \(\{v_m\}_{m=1}^{\infty}\subseteq C^{\infty}(\overline{U})\) 也在 \(W^{1,p}(U)\) 中收敛于 \(u\)。则对任意 \(m\)，
\begin{align*}
\|E_0 u_m-E_0 v_m\|_{W^{1,p}(\mathbb{R}^n)}
\leqslant C''\|u_m-v_m\|_{W^{1,p}(U)}.
\end{align*}
而
\begin{align*}
\|u_m-v_m\|_{W^{1,p}(U)}
\leqslant \|u_m-u\|_{W^{1,p}(U)}+\|u-v_m\|_{W^{1,p}(U)}\to 0\quad (m\to\infty).
\end{align*}
故
\begin{align*}
\|E_0 u_m-E_0 v_m\|_{W^{1,p}(\mathbb{R}^n)}\to 0\quad (m\to\infty).
\end{align*}
由于 \(\{E_0 u_m\}\) 与 \(\{E_0 v_m\}\) 分别收敛于 \(\overline{u}\) 与 \(\overline{v}\)，于是
\begin{align*}
\|\overline{u}-\overline{v}\|_{W^{1,p}(\mathbb{R}^n)}
=\lim\limits_{m\to\infty}\|E_0 u_m-E_0 v_m\|_{W^{1,p}(\mathbb{R}^n)}=0.
\end{align*}
所以 \(\overline{u}=\overline{v}\)。因此 \(E u\) 不依赖于逼近序列的选取。

由 \(E_0\) 的线性性和极限的线性性知\(E\) 是线性的,并且由\eqref{eqJIOFJIOWJIO323u89rfu3}式知
\begin{align*}
\|E u\|_{W^{1,p}(\mathbb{R}^n)}
&=\lim\limits_{m\to\infty}\|E_0 u_m\|_{W^{1,p}(\mathbb{R}^n)}
\leqslant C''\lim\limits_{m\to\infty}\|u_m\|_{W^{1,p}(U)}
=C''\|u\|_{W^{1,p}(U)}.
\end{align*}
故 \(E\) 有界。此外，再证明 \(Eu=u\) a.e. 于 \(U\)。因为 \(u_m\in C^\infty(\overline U)\) 且 \(u_m\to u\) 于 \(W^{1,p}(U)\)，所以 \(u_m\to u\) 于 \(L^p(U)\)。由 \(E_0\) 的定义，对每个 \(m\) 有 \(E_0 u_m=u_m\) a.e. 于 \(U\)，因此
\begin{align*}
E_0 u_m|_U=u_m\to u\quad \text{于 } L^p(U).
\end{align*}
另一方面，由 \(Eu\) 的定义，\(E_0 u_m\to Eu\) 于 \(W^{1,p}(\mathbb R^n)\)，从而限制在 \(U\) 上，有
\begin{align*}
E_0 u_m|_U\to Eu|_U\quad \text{于 } L^p(U).
\end{align*}
由 \(L^p(U)\) 中极限的唯一性，得到
\begin{align*}
Eu|_U=u\quad \text{于 } L^p(U),
\end{align*}
即 \(Eu=u\) a.e. 于 \(U\)。

再证 \(\operatorname{spt}Eu\subseteq V\)。由第6步的构造，所有 \(E_0 u_m\) 的支集都含于紧集
\begin{align*}
K\triangleq \bigcup_{i=0}^N \overline{W_i}.
\end{align*}
因为 \(W_i\subset\subset V\)，所以 \(K\subset\subset V\)。于是对每个 \(m\)，有
\begin{align*}
E_0 u_m=0\quad \text{a.e. 于 } \mathbb R^n\setminus K.
\end{align*}
由 \(E_0 u_m\to Eu\) 于 \(W^{1,p}(\mathbb R^n)\)，限制在 \(\mathbb R^n\setminus K\) 上得到
\begin{align*}
E_0 u_m|_{\mathbb R^n\setminus K}\to Eu|_{\mathbb R^n\setminus K}\quad \text{于 } L^p(\mathbb R^n\setminus K).
\end{align*}
但 \(E_0 u_m|_{\mathbb R^n\setminus K}=0\) 于 \(L^p(\mathbb R^n\setminus K)\)，故
\begin{align*}
Eu|_{\mathbb R^n\setminus K}=0\quad \text{于 } L^p(\mathbb R^n\setminus K),
\end{align*}
即 \(Eu=0\) a.e. 于 \(\mathbb R^n\setminus K\)。因此
\begin{align*}
\operatorname{spt}Eu\subseteq K\subseteq V.
\end{align*}
综上,定理对 \(1\leqslant p<\infty\) 成立。










当 \(p=\infty\) 时，设 \(u\in W^{1,\infty}(U)\)。由 \cref{theorem:ch05-W1infty局部刻画}，\(u\) 在 \(U\) 中局部 Lipschitz 连续。因此，对每个 \(x_0\in\partial U\)，存在开集 \(W\) 包含 \(x_0\)，使得 \(W\subset\subset V\)，且 \(u\) 在 \(W\cap U\) 上是 Lipschitz 连续的，Lipschitz 常数不超过 \(\|Du\|_{L^\infty(U)}\)。

固定这样的 \(W\)。由步骤5中的边界拉直，存在 \(C^1\) 微分同胚 \(\Phi:W\to B\)，\(\Psi=\Phi^{-1}\)，使得
\begin{align*}
\Phi(W\cap U)=B^+,\qquad \Phi(W\cap\partial U)=B\cap\{y_n=0\}.
\end{align*}
定义
\begin{align*}
u'(y)\triangleq u(\Psi(y)),\qquad y\in B^+.
\end{align*}
由于 \(u\) 在 \(W\cap U\) 上 Lipschitz，且 \(\Psi\) 在 \(\overline B\) 上 Lipschitz，所以 \(u'\) 在 \(B^+\) 上是 Lipschitz 的，并且
\begin{align}\label{eq:u'-W1infty估计}
\|u'\|_{W^{1,\infty}(B^+)}
\leqslant C\|u\|_{W^{1,\infty}(W\cap U)}
\leqslant C\|u\|_{W^{1,\infty}(U)}.
\end{align}

对 \(u'\) 应用步骤1--4中 \(p=\infty\) 的反射延拓。具体地，定义
\begin{align*}
\overline u'(y)\triangleq
\begin{cases}
u'(y), & y\in B^+,\\[4pt]
-3u'(y_1,\cdots,y_{n-1},-y_n)
+4u'\left(y_1,\cdots,y_{n-1},-\dfrac{y_n}{2}\right), & y\in B^-.
\end{cases}
\end{align*}
由步骤1--4的论证，\(\overline u'\in W^{1,\infty}(B)\)，满足
\begin{align}\label{eq:延拓u'-W1infty}
\overline u'\big|_{B^+}=u',\qquad
\|\overline u'\|_{W^{1,\infty}(B)}
\leqslant C\|u'\|_{W^{1,\infty}(B^+)}.
\end{align}
令
\begin{align*}
\overline u(x)\triangleq \overline u'(\Phi(x)),\qquad x\in W.
\end{align*}
下面验证 \(\overline u\) 是 \(u\) 到 \(W\) 的延拓，并建立相应的范数估计。

若 \(x\in W\cap U\)，则 \(\Phi(x)\in B^+\)，于是
\begin{align*}
\overline u(x)
=\overline u'(\Phi(x))
=u'(\Phi(x))
=u(\Psi(\Phi(x)))
=u(x).
\end{align*}
因此
\begin{align*}
\overline u\big|_{W\cap U}=u.
\end{align*}

由链式法则和变量替换，类似步骤5(ii)中的估计，结合 \eqref{eq:u'-W1infty估计}、\eqref{eq:延拓u'-W1infty}，得到
\begin{align}\label{eq:局部延拓-W1infty}
\|\overline u\|_{W^{1,\infty}(W)}
\leqslant C\|\overline u'\|_{W^{1,\infty}(B)}
\leqslant C\|u'\|_{W^{1,\infty}(B^+)}
\leqslant C\|u\|_{W^{1,\infty}(W\cap U)}
\leqslant C\|u\|_{W^{1,\infty}(U)}.
\end{align}

由于 \(\partial U\) 紧，存在有限个点 \(x_i^0\in\partial U\)、开集 \(W_i\subset\subset V\) 以及如上构造的局部延拓 \(\overline u_i\)（\(i=1,\cdots,N\)），使得
\begin{align*}
\partial U\subseteq\bigcup_{i=1}^N W_i.
\end{align*}
取 \(W_0\subset\subset U\) 使得
\begin{align*}
U\subseteq\bigcup_{i=0}^N W_i,
\end{align*}
并设 \(\{\zeta_i\}_{i=0}^N\) 为从属于 \(\{W_i\}_{i=0}^N\) 的光滑单位分解。定义
\begin{align*}
\overline u\triangleq\sum_{i=0}^N\zeta_i\overline u_i,\qquad \overline u_0=u.
\end{align*}
由 \eqref{eq:局部延拓-W1infty} 以及单位分解的有限线性组合，得到
\begin{align}\label{eq:整体延拓-W1infty}
\|\overline u\|_{W^{1,\infty}(\mathbb R^n)}
\leqslant C\sum_{i=0}^N\|\overline u_i\|_{W^{1,\infty}(W_i)}
\leqslant C\|u\|_{W^{1,\infty}(U)},
\end{align}
其中常数 \(C\) 只依赖于 \(U,p,n,V\)，不依赖于 \(u\)。又因为 \(W_i\subset\subset V\)，所以
\begin{align*}
\operatorname{spt}\overline u
\subseteq\bigcup_{i=0}^N\operatorname{spt}\zeta_i
\subseteq\bigcup_{i=0}^N W_i
\subset\subset V.
\end{align*}
并且由构造知 \(\overline u|_{U}=u\)。

定义 \(Eu\triangleq\overline u\)。映射 \(u\mapsto Eu\) 是线性的。由 \eqref{eq:整体延拓-W1infty} 知 \(E\) 有界，且满足定理中的三条性质。故定理对 \(p=\infty\) 成立。
\end{proof}

\begin{remark}\label{remark:ch05-延拓注}
\leavevmode
\begin{enumerate}[(i)]
\item 现设 $\partial U$ 为 $C^{2}$。则上面构造的延拓算子 $E$ 也是从 $W^{2,p}(U)$ 到 $W^{2,p}(\mathbb{R}^{n})$ 的有界线性算子。为看到这一点，首先注意在证明的步骤 3、4 中，尽管 $\overline{u}$ 一般不是 $C^{2}$，但它确实属于 $W^{2,p}(B)$。我们还可以由定义\eqref{align:ch05-反射延拓}得到
\begin{align*}
\|\overline{u}\|_{W^{2,p}(B)}\leqslant C\|u\|_{W^{2,p}(B^{+})},
\end{align*}
如前所述，只要 $\partial U$ 为 $C^{2}$，我们由此导出估计
\begin{align}
\|Eu\|_{W^{2,p}(\mathbb{R}^{n})}\leqslant C\|u\|_{W^{2,p}(U)},
\end{align}
其中常数 $C$ 只依赖于 $U$、$V$、$n$ 与 $p$。我们稍后将用到这些观察。
\item 上述构造并未为 $k>2$ 的Sobolev空间 $W^{k,p}(U)$ 提供延拓。这需要更复杂的高阶反射技巧。
\end{enumerate}
\end{remark}

\begin{theorem}[$W^{1,\infty}$的刻画]\label{theorem:ch05-W1infty刻画}
设 $U$是$\mathbb{R}^n$中的有界开集，$\partial U$ 为 $C^{1}$ 类。则 $u:U\to\mathbb{R}$ Lipschitz 连续当且仅当 $u\in W^{1,\infty}(U)$。
\end{theorem}
\begin{remark}
若 $n<p<\infty$，则每个函数 $u\in W^{1,p}$ 都属于 $C^{0,1-\frac{n}{p}}$；但一个以小于 1 的指数H\"older连续的函数未必属于任何Sobolev空间 $W^{1,p}$。
\end{remark}
\begin{proof}

1. 先设 $U=\mathbb{R}^{n}$ 且 $u$ 有紧支集。设 $u\in W^{1,\infty}(\mathbb{R}^{n})$。则 $u_{\varepsilon}\triangleq\eta_{\varepsilon}*u$（$\eta_{\varepsilon}$ 为通常磨光子）光滑，且满足
\begin{align*}
\begin{cases}
u_{\varepsilon}\to u\ \text{一致地，当}\ \varepsilon\to 0,\\
\|Du_{\varepsilon}\|_{L^{\infty}(\mathbb{R}^{n})}\leqslant\|Du\|_{L^{\infty}(\mathbb{R}^{n})}.
\end{cases}
\end{align*}
任取两点 $x,y\in\mathbb{R}^{n}$，$x\neq y$。我们有
\begin{align*}
u_{\varepsilon}(x)-u_{\varepsilon}(y)
=\int_{0}^{1}\frac{\mathrm{d}}{\mathrm{d}t}u_{\varepsilon}(tx+(1-t)y)\,\mathrm{d}t=\int_{0}^{1}Du_{\varepsilon}(tx+(1-t)y)\,\mathrm{d}t\cdot(x-y),
\end{align*}
因而
\begin{align*}
|u_{\varepsilon}(x)-u_{\varepsilon}(y)|
\leqslant\|Du_{\varepsilon}\|_{L^{\infty}(\mathbb{R}^{n})}|x-y|
\leqslant\|Du\|_{L^{\infty}(\mathbb{R}^{n})}|x-y|.
\end{align*}
令 $\varepsilon\to 0$，得
\begin{align*}
|u(x)-u(y)|\leqslant\|Du\|_{L^{\infty}(\mathbb{R}^{n})}|x-y|.
\end{align*}
故 $u$ Lipschitz 连续。

2. 现设 $u:\mathbb{R}^{n}\to\mathbb{R}$ 是 Lipschitz 连续的，记其 Lipschitz 常数为 $\operatorname{Lip}(u)$。固定 $i\in\{1,\cdots,n\}$。对任意 $h\neq 0$，定义差商
\begin{align*}
D_{i}^{-h}u(x)\triangleq\frac{u(x)-u(x-he_{i})}{h}.
\end{align*}
由 Lipschitz 条件，有
\begin{align*}
|D_{i}^{-h}u(x)|\leqslant\operatorname{Lip}(u),\quad \forall x\in\mathbb{R}^{n},
\end{align*}
故
\begin{align*}
\|D_{i}^{-h}u\|_{L^{\infty}(\mathbb{R}^{n})}\leqslant\operatorname{Lip}(u).
\end{align*}
因此差商族 $\{D_{i}^{-h}u\}_{h\neq 0}$ 在 $L^{\infty}(\mathbb{R}^{n})$ 中一致有界，从而在 $L^{2}_{\mathrm{loc}}(\mathbb{R}^{n})$ 中也一致有界。由\cref{Functional Analysis-proposition:Lp对偶与自反性}知$L^{2}_{\mathrm{loc}}(\mathbb{R}^{n})$ 是自反的，再根据\cref{Functional Analysis-proposition:自反空间中弱列紧性等价与有界}, 存在子列$h_{k}\to 0$ 及函数 $v_{i}\in L^{2}_{\mathrm{loc}}(\mathbb{R}^{n})$，使得对每个紧集 $K\subset\mathbb{R}^{n}$，有
\begin{align*}
D_{i}^{-h_{k}}u\rightharpoonup v_{i}\quad\text{在}\ L^{2}(K)\ \text{中弱收敛}.
\end{align*}
对任意 $\phi\in C_c^\infty(\mathbb{R}^n)$，由弱收敛 $D_i^{-h_k}u\rightharpoonup v_i$ 于 $L^2_{\mathrm{loc}}(\mathbb{R}^n)$，有
\begin{align*}
\int_{\mathbb{R}^n}\phi D_i^{-h_k}u\,\mathrm{d}x
\to \int_{\mathbb{R}^n}v_i\,\phi\,\mathrm{d}x.
\end{align*}
另一方面，由 $\|D_i^{-h}u\|_{L^\infty(\mathbb{R}^n)}\leqslant \operatorname{Lip}(u)$ 对一切 $h\neq 0$ 成立，得
\begin{align*}
\left|\int_{\mathbb{R}^n}\phi D_i^{-h_k}u\,\mathrm{d}x\right|
\leqslant \|D_i^{-h_k}u\|_{L^\infty(\mathbb{R}^n)}\int_{\mathbb{R}^n}|\phi|\,\mathrm{d}x
\leqslant \operatorname{Lip}(u)\int_{\mathbb{R}^n}|\phi|\,\mathrm{d}x.
\end{align*}
令 $k\to\infty$，即得
\begin{align*}
\left|\int_{\mathbb{R}^n}v_i\,\phi\,\mathrm{d}x\right|
\leqslant \operatorname{Lip}(u)\int_{\mathbb{R}^n}|\phi|\,\mathrm{d}x.
\end{align*}
由于\cref{Real Analysis-theorem:L无穷空间范数对偶刻画}以及$C_c^\infty(\mathbb{R}^n)$ 在 $L^1(\mathbb{R}^n)$ 中稠密，有
\begin{align*}
\|v_i\|_{L^\infty(\mathbb{R}^n)}
=\sup_{\substack{\phi\in C_c^\infty(\mathbb{R}^n)\\ \|\phi\|_{L^1(\mathbb{R}^n)}= 1}}
\left|\int_{\mathbb{R}^n}v_i\,\phi\,\mathrm{d}x\right|.
\end{align*}
因此综上可得
\begin{align*}
\|v_i\|_{L^\infty(\mathbb{R}^n)}\leqslant \operatorname{Lip}(u).
\end{align*}

任取 $\phi\in C^{\infty}_{c}(\mathbb{R}^{n})$。由\hyperref[theorem:ch05-差商的分部积分公式]{差商的分部积分公式}得
\begin{align*}
\int_{\mathbb{R}^{n}}u(x)D_{i}^{h}\phi(x)\,\mathrm{d}x
=-\int_{\mathbb{R}^{n}}D_{i}^{-h}u(x)\phi(x)\,\mathrm{d}x,
\end{align*}
因为 $\phi$ 光滑且具有紧支集，所以当 $h_{k}\to 0$ 时，$D_{i}^{h_{k}}\phi\to\phi_{x_{i}}$ 在 $L^{\infty}(\mathbb{R}^{n})$ 中成立。于是由\hyperref[Real Analysis-theorem:控制收敛定理]{控制收敛定理}和弱收敛定义得
\begin{align*}
\int_{\mathbb{R}^{n}}u\phi_{x_{i}}\,\mathrm{d}x
=\lim_{k\to\infty}\int_{\mathbb{R}^{n}}uD_{i}^{h_{k}}\phi\,\mathrm{d}x=-\lim_{k\to\infty}\int_{\mathbb{R}^{n}}D_{i}^{-h_{k}}u\,\phi\,\mathrm{d}x
=-\int_{\mathbb{R}^{n}}v_{i}\phi\,\mathrm{d}x.
\end{align*}
上式对一切 $\phi\in C^{\infty}_{c}(\mathbb{R}^{n})$ 成立，故 $v_{i}$ 是 $u$ 关于 $x_{i}$ 的弱导数，即 $v_{i}=u_{x_{i}}$ 在弱意义下成立。由于 $v_{i}\in L^{\infty}(\mathbb{R}^{n})$，所以 $u\in W^{1,\infty}(\mathbb{R}^{n})$。

3. 在 $U$ 有界、$\partial U$ 为 $C^{1}$ 类的一般情形，由\hyperref[theorem:ch05-扩张定理]{扩张定理}，取 $p=\infty$，存在有界线性扩张算子
\begin{align*}
E:W^{1,\infty}(U)\to W^{1,\infty}(\mathbb{R}^{n}),
\end{align*}
使得对每个 $u\in W^{1,\infty}(U)$，有 $Eu=u$ a.e. 于 $U$，且 $Eu$ 在 $\mathbb{R}^{n}$ 中具有紧支集，并且
\begin{align*}
\|Eu\|_{W^{1,\infty}(\mathbb{R}^{n})}\leqslant C\|u\|_{W^{1,\infty}(U)}.
\end{align*}

若 $u\in W^{1,\infty}(U)$，则 $Eu\in W^{1,\infty}(\mathbb{R}^{n})$ 且具有紧支集。由第 1 步的结论，$Eu$ 是 Lipschitz 连续的。因此 $u=Eu|_{U}$ 也是 Lipschitz 连续的。

反之，若 $u:U\to\mathbb{R}$ 是 Lipschitz 连续的，由 \hyperref[theorem:实值-Kirszbraun扩张定理]{Kirszbraun扩张定理}，$u$ 可以扩张为 $\mathbb{R}^{n}$ 上的 Lipschitz 函数 $\overline{u}$，且 $\operatorname{Lip}(\overline{u})=\operatorname{Lip}(u)$。由第 2 步知 $\overline{u}\in W^{1,\infty}(\mathbb{R}^{n})$。限制在 $U$ 上得 $u=\overline{u}|_{U}\in W^{1,\infty}(U)$。

综上，$u$ Lipschitz 连续当且仅当 $u\in W^{1,\infty}(U)$。
\end{proof}



\end{document}
